% DER_XIII-b: Grothendieck's text with the commentary set in, then the overview, inside a group that applies the chunk's own definitions
\dgrpart{1}{XIII-b}{Part 2 of 2, manuscript pages 61--98}{XIII-b\enspace\textperiodcentered\enspace Chapter XIII, part 2 of 2\enspace\textperiodcentered\enspace manuscript pages 61--98}{Chapter XIII \textperiodcentered\ \textup{XIII-b}}
\dgrchunkbegin
% DER_XIII-b: the definitions of the chunk's preamble, applied inside its group
\let\printedpage\relax\newcommand{\printedpage}[1]{\marginpar{\footnotesize p.~#1}}
\let\msp\relax\newcommand{\msp}[1]{\marginpar{\footnotesize [#1]}}
\let\circled\relax\newcommand{\circled}[1]{\tikz[baseline=(c.base)]\node[draw,circle,inner sep=1pt](c){#1};}
\let\cM\relax\newcommand{\cM}{\mathcal{M}}
\let\HO\relax\newcommand{\HO}{\mathrm{HO}}
\let\Hom\relax\newcommand{\Hom}{\mathrm{Hom}}
\let\Ob\relax\newcommand{\Ob}{\mathrm{Ob}}
\let\rf\relax\newcommand{\rf}{\mathrm{f}}
\let\rc\relax\newcommand{\rc}{\mathrm{c}}
\let\M\relax\newcommand{\M}{\mathcal{M}}
\let\Fl\relax\newcommand{\Fl}{\mathrm{Fl}\,}
\let\Filt\relax\newcommand{\Filt}{\mathrm{Filt}}
\let\card\relax\newcommand{\card}{\mathrm{card}\,}
\let\barlongleftrightarrow\relax\newcommand{\barlongleftrightarrow}{\mathrel{\text{\ooalign{$\longleftrightarrow$\cr\hidewidth$\mid$\hidewidth\cr}}}}
\let\idx\relax\newcommand{\idx}[2]{\noindent #1\dotfill #2\par}
\let\entry\relax\newcommand{\entry}[2]{\noindent #1\dotfill\ #2\par}
\let\Mo\relax\newcommand{\Mo}{\mathcal{M}}
\let\Mzg\relax\newcommand{\Mzg}{$\mathbf{M}'_{\mathbf{0}}\mathbf{(g)}$}
\let\Mzd\relax\newcommand{\Mzd}{$\mathbf{M}'_{\mathbf{0}}\mathbf{(d)}$}
\providecommand\thefootnote{}\renewcommand{\thefootnote}{\arabic{footnote}}

% ---- from Der-ChXIII-25-01-04_p036-040.tex
\printedpage{34}
% head: A. Grothendieck : Dérivateurs, Ch. XIII, Édition des Universités de Montpellier II et Paris VII

\begin{itemize}[leftmargin=3em]
\item[c)] \textit{En particulier, si $I'$, $I$ sont des ensembles ordonnés finis, alors tout $u : I' \longrightarrow I$ est admissible. Il en est de même, plus généralement, si $I$, $I'$ sont des ensembles ordonnés et $I'^{\circ}$ est un espace noethérien sobre.}
\end{itemize}

\begin{itemize}[leftmargin=2em]
\item[2)] \textit{Supposons $u : I' \longrightarrow I$ admissible. Considérons}
\[
u_*^W \ : \ \underbrace{\cM(I')_{\rf}W(I')_{\rf}^{-1}}_{\simeq\, \HO_W(I')} \ \longrightarrow \ \underbrace{\cM(I)_{\rf}W(I)_{\rf}^{-1}}_{\HO_W(I)},
\]
\textit{induit par $u_*$. Ce foncteur est adjoint à droite de $u_W^* : \HO_W(I) \longrightarrow \HO_W(I')$.}

\msp{page 61}
\textit{De plus, l'homomorphisme canonique (pour $s \in \Ob I$, $X' \in \Ob \HO_W(I')$)}
\[
u_*^W(X')(s) \ \simeq \ p_{s\,*}(X'|I'/s) \ \overset{\text{déf}}{=} \ \underset{I'/s}{\underleftarrow{\mathrm{holim}}}{}_W(X'|I'/s)
\]
\textit{(où $p_s : I'/s \longrightarrow e$) est un isomorphisme, si on suppose que les $I'/s$ sont admissibles (on dira alors que $u$ est} strictement \textit{admissible).}
\end{itemize}

\textbf{Corollaire.} \textit{Sur le domaine formé des ensembles ordonnés finis, le prédérivateur} $\HO_W$ \textit{satisfait (en plus des axiomes} \textbf{Der}~1, \textbf{Der}~2\textit{) aux axiomes} \textbf{Der}~4 \textit{et} \textbf{Der}~5 \textit{(existence des images directes $u_*$, adjoints à droite des foncteurs $u^*$, et calcul standard des fibres d'une image directe).}

\textsc{Démonstration} du théorème-scholie.

$1^\circ$) \textit{a}) est immédiat. Pour \textit{b}), on procède par récurrence sur le cardinal $n$ de l'ensemble ordonné fini $I$, les cas $n \in \{0,1\}$ étant triviaux. Supposons donc $n \geq 2$, et le théorème prouvé pour les $n' < n$. Je distingue deux cas : si $I$ admet un plus grand élément $a$, alors, désignant par $\alpha : e \longrightarrow I$ le foncteur de valeur $a$, par $p : I \longrightarrow e$ la projection, on sait que l'on a
\[
p_* \ \simeq \ \alpha^*, \qquad \textit{i.e.} \ \varprojlim_{I} X_i \simeq X_a,
\]
d'où le fait que sur $\cM(I)$ tout entier
\msp{page 62}
(pas seulement sur $\cM(I)_{\rf}$), $p_*$ est compatible aux localiseurs $W(I)$, $W$.

\dgrblock{dgrb716}{Finite limits assembled from a two-piece cover}
\begin{dcomment}{exposition}{p61}{Finite limits assembled from a two-piece cover}
Grothendieck's finite-poset induction makes the geometry of a homotopy limit visible. Instead of treating a large compatibility diagram at once, he cuts the indexing poset into two lower sets and glues the two resulting limits over their overlap. This argument remains useful even though the Brown lemma gives a shorter proof of admissibility under the stronger standing hypotheses.

Let $I=I_1\cup I_2$ with $I_1,I_2$ lower subsets, and put $I_0=I_1\cap I_2$. Define $v:I\to\Psi=(1\leftarrow0\to2)$ by sending the overlap to $0$, the points belonging only to $I_1$ to $1$, and those belonging only to $I_2$ to $2$. The lower-set condition makes this assignment order preserving: an arrow cannot travel from one exclusive part into the other. Its comma categories over $0,1,2$ are precisely $I_0,I_1,I_2$. Right Kan extension along $v$ therefore produces the span
\[
 \varprojlim_{I_1}X\longrightarrow\varprojlim_{I_0}X
 \longleftarrow\varprojlim_{I_2}X.
\]
Taking its pullback gives $\varprojlim_I X$. Indeed, a cone to $X$ is exactly a pair of cones to its restrictions whose two cones on $I_0$ agree. This proves the ordinary limit identity by its universal property, including the cases in which the overlap is empty.

For a finite nonempty poset without a greatest element, partition its maximal elements into two nonempty sets $M_1,M_2$ and take $I_j=\bigcup_{m\in M_j}\{i:i\leq m\}$. Every point lies below a maximal element, so the lower sets cover $I$, and neither is all of $I$. Their overlap is also smaller. The inductive argument is therefore well founded. If $I$ has a greatest element $a$, the limit is $X_a$ in the source's contravariant convention; the empty poset has the terminal object as its limit.

The homotopical step is to ensure that the span just constructed is fibrant. Right Kan extension preserves fibrancy, and a fibrant $\Psi$-diagram consists of fibrant objects $A_0,A_1,A_2$ with both arrows $A_j\to A_0$ fibrations. Its ordinary pullback then represents its homotopy pullback. The induction replaces ``take an arbitrary pullback'' by ``take a pullback whose legs satisfy a lifting condition''. The latter is the operation stable under weak equivalences of the input span.

Composition of admissible functors is also transparent at this point. If $u$ and $v$ preserve weak equivalences between fibrant diagrams after right Kan extension, then $u_*$ first yields another fibrant diagram, to which $v_*$ can be applied. The canonical identification $(vu)_*\cong v_*u_*$ proves admissibility of the composite. No new comparison of replacements is required. This explains why the elementary shape $\Psi$ controls the whole finite induction.
\end{dcomment}


Si $I$ n'a pas de plus grand élément, alors l'ensemble $M$ des ses éléments maximaux est de cardinal $\geq 2$, on peut donc l'écrire $M = M_1 \amalg M_2$, et on aura % [sic] des ses -> de ses, typographical slip
\[
I \ = \ I_1 \cup I_2, \qquad \text{où } I_1 = I/M_1,\ I_2 = I/M_2
\]
sont deux ouverts de $I$, différents de $I$, donc de cardinaux $n_1$, $n_2 \ < \ n$. Considérons aussi $I_0 = I_1 \cap I_2$, c'est un ouvert de cardinal $< n$. Le diagramme
\[
\begin{tikzcd}[row sep=small]
 & I_1 \arrow[dr, dotted] & \\
I_0 \arrow[ur] \arrow[dr] & & I = I_1 \cup I_2 \\
 & I_2 \arrow[ur, dotted] &
\end{tikzcd}
\]
% foot: Version 25 janvier 2004
\printedpage{35}
% head: A. Grothendieck : Dérivateurs, Ch. XIII, Édition des Universités de Montpellier II et Paris VII
correspond à une flèche
\[
\begin{tikzcd}[row sep=small, column sep=small]
 & & & 1 & \\
u \ : \ I \arrow[r] & \Psi = 0 \arrow[urr, start anchor=east, end anchor=south west] \arrow[drr, start anchor=east, end anchor=north west] & & & , \\
 & & & 2 &
\end{tikzcd}
\]
avec
\[
I_1 \ = \ I/1, \qquad I_2 \ = \ I/2, \qquad I_0 \ = \ I/0.
\]
L'hypothèse de récurrence et \textit{a}) impliquent que $u$ est admissible. Comme un composé de foncteurs admissibles est admissible (tautologiquement, via $(vu)_* \simeq v_*u_*$), on est réduit finalement au
\msp{page 63}

\textbf{Lemme.} \textit{La catégorie} \begin{tikzcd}[row sep=small, column sep=small, baseline=(\tikzcdmatrixname-2-1.base)]
 & 1 \\
\Psi = 0 \arrow[ur] \arrow[dr] & \\
 & 2
\end{tikzcd} \quad \textit{est admissible.}

Ceci va résulter du

\textbf{Corollaire} (\footnote[36]{\textbf{NB} Il y a une variante évidente, caractérisant les $X$ tels que $X \longrightarrow e$ soit dans $TF(\Psi)$ comme ceux pour lesquels $X_0$, $X_1$, $X_2$ sont fibrants \textit{triviaux,} et $\alpha_1$, $\alpha_2 \in TF$.}). \textit{Soit} \begin{tikzcd}[row sep=small, column sep=small, baseline=(\tikzcdmatrixname-2-1.base)]
 & X_1 \arrow[dl, "\alpha_1"'] \\
X_\bullet = X_0 & \\
 & X_2 \arrow[ul, "\alpha_2"]
\end{tikzcd} \quad \textit{dans $\cM(\Psi)$. Pour que $X_\bullet$ soit fibrant, il faut et il suffit que l'on ait}

\begin{itemize}[leftmargin=3em]
\item[a)] \textit{$X_0$, $X_1$, $X_2$ sont fibrants, et}
\item[b)] \textit{$\alpha_1$, $\alpha_2 \in F$.}
\end{itemize}

La nécessité de \textit{a}) est cas particulier du corollaire 2, page 48, celle de \textit{b}) est cas particulier du corollaire 1 page 50. Pour la suffisance (dont on n'aura pas besoin pour le théorème), on sait montrer que dans le diagramme commutatif à traits pleins
\[
% Audit resolved: PDF leaf 7 confirms interrupted vertical sides A_0 -> X_0 and A_2 -> X_2, obscured by the displaced B labels; the printed geometry is retained. See corrections.csv.
\begin{tikzcd}
A_1 \arrow[r] \arrow[dr, "TC"] \arrow[dd] & A_0 \arrow[dr, "TC"] \arrow[d, dash] & A_2 \arrow[l] \arrow[dr, "TC"] \arrow[d, dash] & \\
 & B_1 \arrow[r] \arrow[dl, dotted, "h_1"'] \arrow[d] & B_0 \arrow[dl, dotted, "h_0"'] \arrow[d] & B_2 \arrow[l] \arrow[dl, dotted, "h_2"'] \\
X_1 \arrow[r, "\alpha_1"'] & X_0 & X_2 \ , \arrow[l, "\alpha_2"] &
\end{tikzcd}
\]
avec trois flèches marquées $TC$, on peut trouver $h_1$, $h_0$, $h_2$ laissant le diagramme commutatif (trois triangles et deux carrés commutatifs à respecter). Pour ceci, on commence par choisir $h_0$ de façon à rendre le triangle vertical médian commutatif -- c'est possible car $X_0$ est fibrant, $A_0 \longrightarrow B_0$ dans $TC$.

% foot: Version 25 janvier 2004
\printedpage{36}
% head: A. Grothendieck : Dérivateurs, Ch. XIII, Édition des Universités de Montpellier II et Paris VII
\msp{page 64}

\dgrblock{dgrb717}{Matching objects explain the lifting diagram}
\begin{dcomment}{examples}{p63}{Matching objects explain the lifting diagram}
Grothendieck's three lifts express a recursive compatibility condition. The lift at $X_0$ must be chosen first, because the two remaining lifts are required to induce that same map after composition with $X_1\to X_0$ and $X_2\to X_0$. Choosing three unrelated lifts from objectwise fibrancy would not enforce the squares. The fibrations on the two legs are exactly what turns each final choice into a lifting problem with a prescribed target.

The calculation is the smallest instance of the matching-object description for an inverse diagram. Let $I$ be a finite poset, still indexing contravariant diagrams. For $i\in I$ put
\[
                M_iX=\varprojlim_{j<i}X_j.
\]
An injectively fibrant diagram is characterized by the requirement that each natural map $X_i\to M_iX$ be a fibration. When $i$ has no smaller element, the matching object is terminal, so this says that $X_i$ is fibrant. For $I=\Psi$, the matching objects are $M_0X=e$ and $M_1X=M_2X=X_0$, giving exactly the corollary. This is the inverse-category case of the familiar Reedy description \citep[Theorem 5.1.3]{Hovey1999}.

The lifting criterion can be proved without invoking the terminology. Choose a linear ordering of the finite set of vertices in which smaller vertices occur first. Given an objectwise trivial cofibration $A\to B$ and a map $A\to X$, suppose the lift $B_j\to X_j$ has been chosen for every $j<i$. These maps induce $B_i\to M_iX$. The desired map $B_i\to X_i$ is now a lift in the square whose left side is $A_i\to B_i$ and whose right side is $X_i\to M_iX$. Proceeding through the vertices gives a natural transformation $B\to X$.

For necessity, start with a lifting problem $A\to X_i$, $B\to M_iX$, where $A\to B$ is a trivial cofibration. Form two diagrams which equal $X_j$ for $j<i$, take the values $A$ and $B$ at $i$, respectively, and equal the initial object outside the lower subset $\{j:j\leq i\}$. The given maps to the matching object specify their transition maps. Extension by the initial object is consistent because this subset is lower. Their comparison is objectwise a trivial cofibration, and the first diagram maps to $X$ by the identity below $i$ and the given map at $i$. A lift against $X\to e$ supplies the required $B\to X_i$. This proves the converse. The same finite recursion, factoring $X_i\to M_i(RX)$ at each stage as a trivial cofibration followed by a fibration, constructs a fibrant replacement $X\to RX$.

The condition becomes more informative at a branching vertex. Consider the diamond $0<1<3$, $0<2<3$. Its top matching object is
\[
                           M_3X=X_1\times_{X_0}X_2.
\]
In cochain complexes of vector spaces over a field $k$, take $X_0=0$ and $X_1=X_2=X_3=k[0]$, with both maps from $X_3$ equal to the identity. All objects are fibrant in the usual model structure, and every nonidentity arrow of this diagram is a fibration: the maps are degreewise surjective. Yet
\[
                  X_3\longrightarrow M_3X=k[0]\oplus k[0]
\]
is the diagonal and is not surjective. The diagram is not injectively fibrant. Thus the two separate legs happen to suffice for $\Psi$, whereas a general poset requires compatibility with the entire lower matching diagram. This distinction is why replacements of diagrams carry more information than separate replacements of their components.
\end{dcomment}


Ensuite, on voit séparément pour $h_1$, $h_2$, qu'on peut les trouver en satisfaisant chacune aux deux contraintes de commutativité (1 triangle vertical et un carré penché). Pour $h_1$ p.~ex. l'existence se voit sur le diagramme carré auxiliaire (traits pleins)
\[
\begin{tikzcd}
A_1 \arrow[r, "TC"] \arrow[dd] & B_1 \arrow[dd] \arrow[ddl, dotted, "h_1"'] \arrow[dr, dotted, "\alpha_1^B"] & \\
 & & B_0 \arrow[dl, dotted, "h_0"] \\
X_1 \arrow[r, "F", "\alpha_1^X"'] & X_0 \ \ . &
\end{tikzcd}
\]

\textbf{NB} C'est la démonstration de la proposition 3, page 50.

Le corollaire (partie “il faut”) étant vu, le lemme en résulte, car on sait que si on a un homomorphisme $f$ de carrés cartésiens $X$, $Y$
\[
\begin{tikzcd}[row sep=small, column sep=small]
 & X_1 \arrow[dl, "p_1^X"'] \arrow[dd, "f_1"'] & & X_3 \arrow[ll] \arrow[dl] \arrow[dd, "f_3"] \\
X_0 \arrow[dd, "f_0"'] & & X_2 \arrow[ll, crossing over, "p_2^X"'] & \\
 & Y_1 \arrow[dl, "p_1^Y"'] & & Y_3 \arrow[ll] \arrow[dl] \\
Y_0 & & Y_2 \ \ , \arrow[ll, "p_2^Y"'] \arrow[from=uu, crossing over, "f_2"] &
\end{tikzcd}
\]
avec
\[
p_1^X,\ p_2^X,\ p_1^Y,\ p_2^Y \ \ \text{fibrants \texttt{[appartenant à} $F$\texttt{]} (il suffit même $p_1^X$ et $p_1^Y$ fibrant),}
\]
et
\[
f_0,\ f_1,\ f_2 \ \in \ W,
\]
alors $f_3 \in W$. (Ça a été vu dans la théorie des produits fibrés $W$-homotopiques \texttt{[voir chapitre XII, p. 253]}, et la vérification
\msp{page 65}
directe est immédiate.) De l'hypothèse initiale pertinente que \begin{tikzcd}[row sep=small, column sep=small, baseline=(\tikzcdmatrixname-2-1.base)]
 & X_1 \arrow[dl] \\
X_0 & \\
 & X_2 \arrow[ul]
\end{tikzcd} et \begin{tikzcd}[row sep=small, column sep=small, baseline=(\tikzcdmatrixname-2-1.base)]
 & Y_1 \arrow[dl] \\
Y_0 & \\
 & Y_2 \arrow[ul]
\end{tikzcd}
sont fibrants, on n'a besoin d'utiliser qu'un petit bout, p.~ex. que $p_1^X$ et $p_1^Y$ sont fibrants, qui assure que les deux carrés $X$, $Y$ sont $W$-homotopiquement cartésiens, ce qui suffit à impliquer alors que $f_0$, $f_1$, $f_2 \in W \Longrightarrow f_3 \in W$.

\dgrblock{dgrb718}{A homotopy pullback computed as a complex}
\begin{dcomment}{examples}{p64}{A homotopy pullback computed as a complex}
Grothendieck uses the pullback lemma to make finite gluing invariant under weak equivalences. A concrete calculation shows what fibrant replacement adds to an ordinary pullback. Work with cochain complexes over a field and consider maps $a:X_1\to X_0$ and $b:X_2\to X_0$. A model for their homotopy pullback is
\[
 H=\operatorname{Cone}(a-b:X_1\oplus X_2\to X_0)[-1].
\]
With coordinates chosen so that
\[
 H^n=X_1^n\oplus X_2^n\oplus X_0^{n-1},
 \quad d(x_1,x_2,h)=(dx_1,dx_2,-dh+a(x_1)-b(x_2)),
\]
the equation $d^2=0$ follows from the fact that $a$ and $b$ are chain maps. The third coordinate records the homotopy witnessing equality of the first two images. A strict pullback asks those images to be equal without this additional coordinate.

To verify that this complex computes the derived pullback, replace $a$ by the surjective map $\widetilde X_1\to X_0$, where
\[
 \widetilde X_1^n=X_1^n\oplus X_0^n\oplus X_0^{n-1},
 \quad d(x,z,h)=(dx,dz,-dh+a(x)-z),
\]
and the map to $X_0$ is $(x,z,h)\mapsto z$. The inclusion $x\mapsto(x,a(x),0)$ is a quasi-isomorphism: the projection back to $X_1$ has kernel the contractible complex with differential $(z,h)\mapsto(dz,-dh-z)$, contracted by $(z,h)\mapsto(-h,0)$. The new map is degreewise surjective and hence a fibration. Its ordinary pullback along $b$ imposes $z=b(x_2)$ and is exactly the displayed complex $H$.

There is a short exact sequence of complexes
\[
 0\longrightarrow X_0[-1]\longrightarrow H
   \longrightarrow X_1\oplus X_2\longrightarrow0.
\]
Its connecting map on cohomology is induced by $a-b$. This gives the usual exact sequence for a homotopy pullback and also proves invariance directly: a componentwise quasi-isomorphism between two such spans induces an isomorphism on the two outside terms of the associated long exact sequences, and hence on the middle term. If $a-b$ is degreewise surjective, its ordinary kernel maps quasi-isomorphically to $H$. In particular this occurs when either $a$ or $b$ is degreewise surjective, exactly the fibration condition relevant to this model.

For the span $0\to k[0]\leftarrow0$, the ordinary pullback is zero, while $H=k[-1]$ has $H^1(H)=k$. A fibration replacement of one leg restores this missing class. This is the algebraic version of the familiar fact that the homotopy pullback of the basepoint maps $*\to Y\leftarrow *$ is the loop object of $Y$. It explains both why the manuscript requires fibrant legs and why its gluing step cannot be replaced by taking limits in the homotopy category with no coherent diagram behind them.

In the source's general framework, the corresponding invariance statement uses the $W$-fibration property included in $\mathbf M''$, or the equivalent derived-functor argument under the closed exact hypotheses. No failure of properness is being overlooked in the proof on pp.~64--65: the relevant base-change compatibility is already among its assumptions. The statement that one fibrant leg suffices concerns this particular pullback computation, not a claim that all transition maps in an arbitrary diagram can be ignored.
\end{dcomment}


Je vais essayer de généraliser ce raisonnement, pour traiter le cas d'un ensemble ordonné $I$ qui peut être infini, mais tel que $I^\circ$ soit noethérien et sobre, donc toute suite croissante de parties fermées de $I$, et toute suite décroissante de parties ouvertes, est stationnaire, et de plus toute partie ouverte filtrante croissante a un plus grand élément.

% foot: Version 25 janvier 2004
\printedpage{37}
% head: A. Grothendieck : Dérivateurs, Ch. XIII, Édition des Universités de Montpellier II et Paris VII

% margin: a vertical bar is printed in the left margin beside the lines "prouver ceci : (*) Si I est tel que ... alors I aussi."
On veut prouver que $I$ est admissible, et par “récurrence noethérienne” on est ramené à prouver ceci : $(*)$ Si $I$ est tel que pour tout ouvert $I' \subset I$, avec $I' \neq I$, $I'$ soit admissible, alors $I$ aussi.

$(*)$ admis, prouvons que $I = I_0$ est admissible.

\msp{page 66}
S'il ne l'était, par $(*)$ il existerait $I_1 \subset I_0$, $I_1 \neq I_0$, tel que $I_1$ non admissible. De même, il y aurait $I_2 \subset I_1$, $I_2 \neq I_1$, non admissible etc. -- on aurait construit une suite strictement décroissante d'ouverts $I_n$ de $I$, contrairement au caractère noethérien de $I^\circ$.

Prouvons $(*)$, et distinguons encore deux cas.

\textit{a}) $I^\circ$ est irréductible, donc admet un point générique, \textit{i.e.} un plus petit élément, \textit{i.e.} $I$ admet un plus grand élément. Ce cas a déjà été traité tantôt -- toute catégorie avec objet final est admissible.

\textit{b}) $I^\circ$ est réductible, donc réunion de deux fermés $I_1^\circ$, $I_2^\circ$ distincts de $I^\circ$, donc $I$ est réunion de deux ouverts $I_1$, $I_2$ distincts de $I$. Alors par l'hypothèse dans $(*)$, $I_1$, $I_2$ et $I_0 = I_1 \cap I_2$ sont admissibles. On termine alors comme tantôt, étant ramené à l'admissibilité de \begin{tikzcd}[row sep=small, column sep=small, baseline=(\tikzcdmatrixname-2-1.base)]
 & 1 & \\
\Psi = 0 \arrow[ur] \arrow[dr] & & ! \\
 & 2 &
\end{tikzcd}

\msp{page 67}
Prouvons la partie \textit{c}) du théorème-scholie, partie $1^\circ$). Cela résulte de \textit{a}), une fois vu que pour $s \in I$, $I'/s$ satisfait aux mêmes hypothèses que $I'$. Or $I$ étant ordonné, $I/s$ est une sous-catégorie ouverte de $I$, donc $I'/s$ une sous-catégorie ouverte de $I'$, dont $(I'/s)^\circ$ un sous-ensemble ordonné fermé de $I'^{\circ}$. Ce dernier étant noethérien et sobre, il en est de même de tout fermé, d'où la conclusion.

\dgrblock{dgrb719}{The proposed noetherian enlargement is actually finite}
\begin{dcomment}{exposition}{p66}{The proposed noetherian enlargement is actually finite}
Grothendieck's noetherian induction is patterned on the decomposition of a closed space into irreducible pieces. A generic point makes one piece have a greatest indexing object, and a reducible piece can be glued from two smaller ones. The induction is valid. For the specific Alexandrov topology of a poset, however, the combination ``noetherian and sober'' is stronger than it first appears: it forces the poset to be finite. Thus this paragraph does not in fact add infinite diagram shapes.

We verify the assertion in the precise order convention used here. Opens of $I$ are lower subsets, so opens of $I^{\mathrm{op}}$ are upper subsets of $I$. Sobriety of $I^{\mathrm{op}}$ says that every nonempty directed lower subset of $I$ has a greatest element. If $x_0<x_1<\cdots$ were a strictly increasing sequence, its downward closure would be directed. It could not have a greatest element: any alleged greatest element belongs below some $x_n$, whereas $x_{n+1}$ is strictly larger. Consequently every nonempty subset of $I$ has a maximal element; otherwise dependent choice would construct such a sequence. This is the order-theoretic sobriety criterion, with the order reversed, in \citet[Proposition B.5]{Lindenhovius2014}.

Let $X$ denote $I$ with its upper topology, that is, the topology of $I^{\mathrm{op}}$. Every nonempty subspace $S\subset X$ has an isolated point: take a maximal element $m$ of $S$, so that $S\cap\{x:m\leq x\}=\{m\}$. We now use noetherianity. The union $U$ of all finite open subsets of $X$ is open. Every open subset of a noetherian space is quasi-compact, so the cover of $U$ by its finite open subsets has a finite subcover. Hence $U$ itself is finite. If $X\setminus U$ were nonempty, choose an isolated point $x$ of that subspace. There would be an open $V\subset X$ with $V\cap(X\setminus U)=\{x\}$. Then $U\cup V=U\cup\{x\}$ would be a finite open subset not contained in $U$, contradicting the definition of $U$. Therefore $X=U$ is finite.

Conversely, every finite poset has a noetherian sober Alexandrov topology. Noetherianity is immediate. A nonempty irreducible closed subset has finitely many minimal elements in the lower-topology convention, and irreducibility forces exactly one of them; the closed subset is its closure. Thus the two hypotheses in the proposed enlargement are equivalent to finiteness for posets.

This observation should not be confused with the false assertion that every noetherian sober topological space is finite. Spectra supply abundant infinite examples. The essential extra hypothesis is the Alexandrov topology, in which arbitrary intersections of opens remain open. Nor is it necessary to insert a defect into the manuscript's induction: its irreducible and reducible cases work. The correction is to its claimed extra scope. Infinite towers and group actions on the following pages require another argument, which the Brown lemma supplies whenever the necessary diagram structures and limits exist.
\end{dcomment}


Prouvons la partie $2^\circ$) du théorème-scholie. Considérons
\[
\begin{tikzcd}[row sep=large]
\cM(I')_{\rc} \arrow[r, hook, dotted] & \cM(I') & \cM(I')_{\rf} \arrow[l, hook', dotted] \arrow[d, "u_*^{\rf}"'] \\
\cM(I)_{\rc} \arrow[r, hook, dotted] \arrow[u, "u_{\rc}^*"] & \cM(I) \arrow[u, "u^*"] & \cM(I)_{\rf} \ , \arrow[l, hook', dotted]
\end{tikzcd}
\]
où $u^*$ et a fortiori $u_{\rc}^*$ sont compatibles avec $W(I)$, $W(I')$, et par hypothèse $u_*$ est compatible avec $W(I')_{\rf}$, $W(I)_{\rf}$. En passant aux catégories de fractions, on trouve

% foot: Version 25 janvier 2004
\printedpage{38}
% head: A. Grothendieck : Dérivateurs, Ch. XIII, Édition des Universités de Montpellier II et Paris VII
\msp{page 68}
\[
\begin{tikzcd}[row sep=large]
\cM(I')_{\rc}W(I')_{\rc}^{-1} \arrow[r, dotted, "\sim"] & \HO_W(I') & \cM(I')_{\rf} \arrow[l, dotted, "\sim"'] \arrow[d, "u_*^W"'] \\
\cM(I)_{\rc}W(I)_{\rc}^{-1} \arrow[r, dotted, "\sim"] \arrow[u, "u_W^*"] & \HO_W(I) & \cM(I)_{\rf} \arrow[l, dotted, "\sim"']
\end{tikzcd}
\]
\texttt{[plutôt} $\cM(I')_{\rf}W(I')_{\rf}^{-1}$ \texttt{(resp.} $\cM(I)_{\rf}W(I)_{\rf}^{-1}$\texttt{) au lieu de} $\cM(I')_{\rf}$ \texttt{(resp.} $\cM(I)_{\rf}$\texttt{)]} et on veut montrer que $u_W^*$ et $u_*^W$ sont adjoints, de façon précise trouver un isomorphisme bifonctoriel, pour $F$ \textit{dans} $\cM(I)_{\rc}$ et $G'$ \textit{dans} $\cM(I')_{\rf}$
\[
\Hom_W(u^*F, G') \xrightarrow{\ \sim\ } \Hom_W(F, u_*G') \tag*{$(*)$}
\]
(où $\Hom_W$ désigne le Hom dans $\HO_W(I')$ resp. $\HO_W(I)$). Comme $F$ est dans $\cM(I)_{\rc}$, $u^*F$ est dans $\cM(I')_{\rc}$, et comme $G'$ est dans $\cM(I')_{\rf}$, $u_*G'$ est dans $\cM(I)_{\rf}$. Il s'ensuit que les deux membres de $(*)$ sont des quotients des deux membres dans la bijection canonique
\[
\Hom_{\cM(I')}(u^*F, G') \xrightarrow[\ \sim\ ]{\alpha} \Hom_{\cM(I)}(F, u_*G'). \tag*{$(**)$}
\]
Il faut simplement montrer que cette dernière transforme la relation d'équivalence pertinente $R'$ du premier membre,
\msp{page 69}
en celle $R$ du second.

\dgrblock{dgrb720}{The adjunction after localization, with both inverses}
\begin{dcomment}{exposition}{p68}{The adjunction after localization, with both inverses}
Grothendieck wants the limit functor to retain its universal property after replacement. Knowing that it sends weak equivalences to weak equivalences makes it descend to the homotopy category; it does not alone identify its adjoint. The argument on pp.~68--71 checks that the ordinary adjunction identifies the two homotopy relations. We can make its two directions explicit without choosing a globally functorial replacement.

Choose $F$ cofibrant in $\mathcal M(I)$ and $G$ fibrant in $\mathcal M(I')$. Then $u^*F$ is cofibrant and $u_*G$ is fibrant. Every morphism between these objects in the respective homotopy categories is represented by an ordinary map. Restriction followed by the ordinary counit defines
\[
 \beta:\operatorname{Hom}_{\mathrm{HO}(I)}(F,u_*G)
 \longrightarrow\operatorname{Hom}_{\mathrm{HO}(I')}(u^*F,G).
\]
It is well defined because restriction preserves all weak equivalences. It is surjective because every ordinary map $u^*F\to G$ has a strict adjoint $F\to u_*G$. This is the one direction available before admissibility has been used.

To construct the inverse, choose a trivial cofibration $i':u^*F\to F'$ with $F'$ fibrant, and let $i:F\to u_*F'$ be its strict adjoint. The object $F'$ is also cofibrant. Composition with $i'$ induces a bijection
\[
 \operatorname{Hom}_{\mathrm{HO}(I')}(F',G)
    \xrightarrow{\sim}\operatorname{Hom}_{\mathrm{HO}(I')}(u^*F,G).
\]
For a class $g$ on the right, let $a$ be its unique inverse image and define
\[
                            \gamma(g)=u_*^W(a)\circ[i].
\]
This is where admissibility enters: $u_*^W$ is a functor on the localization of fibrant objects, so it acts on the class $a$, independently of its representative. Naturality of the ordinary counit, together with the fact that $i'$ is adjoint to $i$, gives $\beta\gamma(g)=a[i']=g$.

For the other composite, represent a class $b:F\to u_*G$ by an ordinary map and let $g:u^*F\to G$ be its ordinary adjoint. Since $i'$ is a trivial cofibration and $G$ is fibrant, $g$ extends to an ordinary map $a:F'\to G$ satisfying $ai'=g$. Its class is the inverse image used in defining $\gamma$. The ordinary adjunction then gives $u_*a\circ i=b$, proving $\gamma\beta([b])=[b]$. Thus $\beta$ is bijective. Its naturality, already built into restriction and the counit, gives the required bifunctorial adjunction.

This is the derived-adjunction construction; compare \citet[Lemma 1.3.10]{Hovey1999}. Independence of the chosen $F'$ follows from uniqueness of the inverse of $\beta$, rather than from an additional comparison of arbitrary replacements. It also explains the manuscript's remark that functoriality in $G$ appears a posteriori. The adjoint obtained is
\[
                                u_W^*\dashv u_*^W.
\]
The two occurrences of ``adjoint \`a gauche'' on p.~71 are slips in the source and must read ``adjoint \`a droite''. The Hom formula itself already fixes the orientation, so the proof is unaffected. Likewise the target containing $u^*F'$ on p.~70 is ill typed: $F'$ is already an $I'$-diagram, and the intended object is $u^*F$.
\end{dcomment}


Je dis que l'on a de toute façon
\[
R \ \leq \ \alpha(R')
\]
\textit{sans} hypothèse d'admissibilité sur $u$, \textit{i.e.} que la bijection invese de $(**)$ donne par passage aux quotients un homomorphisme % [sic] invese -> inverse, missing letter
\[
\Hom_W(u^*F, G') \longleftarrow \Hom_W(F, u_*G'), \tag*{$(***)$}
\]
forcément surjectif, et de plus fonctoriel en $F$ \texttt{[,} $G'$ \texttt{barré]}, variables dans $\HO_W(I)$ et $\HO_W(I')$ respectivement. En effet, on définit la flèche $(*$ bis$)$ \texttt{[plutôt} $(***)$\texttt{]} comme composé
\[
\begin{tikzcd}[column sep=large]
\Hom_W(F, u_*G') \arrow[r, "\underbrace{\hspace{2.6cm}}_{\text{\footnotesize\begin{tabular}{@{}c@{}}car $u^*$ est\\ compatible avec\\ $W(I)$, $W(I')$ sur\\ $\cM(I)$ tout entier,\\ et on interprète\\ les deux membres\\ comme des Hom\\ dans $\cM(I)W(I)^{-1}$\\ et $\cM(I')W(I')^{-1}$.\end{tabular}}}"'] & \Hom_W(u^*F, u^*u_*G') \arrow[dr, "\text{\footnotesize\begin{tabular}{@{}l@{}}composition\\ avec l'homomorphisme\\ d'adjonction $u^*u_*G' \longrightarrow G'$\end{tabular}}"] & \\
 & & \Hom_W(u^*F, G').
\end{tikzcd}
\]
% foot: Version 25 janvier 2004

% ---- from Der-ChXIII-25-01-04_p041-045.tex
\printedpage{39}
% head: A. Grothendieck : Dérivateurs, Ch. XIII, Édition des Universités de Montpellier II et Paris VII
% foot: Version 25 janvier 2004

Cela implique donc que le foncteur sur $\HO_W(I)$ $F \longmapsto \Hom_{\HO_{W}(I')}(u^*F,G')$ est, sinon représentable, du moins quotient terme à terme du foncteur représenté par $u_*G'$ (\footnote[37]{Mais il n'est pas clair, sans hypothèse supplémentaire sur la situation, que ce foncteur représentable puisse être choisi en dépendance fonctorielle de $G'$, envisagé dans $\HO_W(I')$.}).

\msp{page 70}

Mais supposons maintenant que $u$ soit admissible. Alors je vais définir la flèche en sens initial $(*)$ (p. 68) de la façon suivante : Je choisis une flèche
\[
i' \ : \ u^*F \longrightarrow F' \qquad \text{avec } i' \in TC \subset C \cap W,\ F' \in \Ob \M(I')_{\rf}
\]
(donc en fait $F' \in \Ob \M(I')_{\mathrm{cf}}$, car $u^*F$ est cofibrant), et par adjonction la flèche $i'$ correspond à
\[
i \ : \ F \longrightarrow u_*F'.
\]
Je considère le diagramme (commutatif!)
\[
\begin{tikzcd}[column sep=large, row sep=large]
\Hom(F',G') \arrow[r, "\alpha"] \arrow[d, "{\scriptstyle \lambda:f'\,\longmapsto\, f'\circ i'}"] & \Hom(u_*F',u_*G') \arrow[d, "{\scriptstyle \mu:f\,\longmapsto\, f\circ i}"] \\
\Hom(u^*F,G') \arrow[r, "\text{adj.}"] & \Hom(F,u_*G')\ .
\end{tikzcd}
\]

On note que l'on a encore $F'$ cofibrant (car $u^*F$ l'est et $i' \in TC \subset C$), donc $\Hom_W(F',G')$ est un quotient de $\Hom(F',G')$, et $\lambda$ induit un isomorphisme
\[
\Hom_W(F',G') \xrightarrow{\ \sim\ } \Hom_W(u^*F',G') % [sic] u^*F' -> u^*F, the target of lambda is Hom(u^*F,G')
\]
(car $i' \in TC \subset W$). Comme $F'$, $G' \in \Ob \M(I')_{\rf}$ et que $u_*^{\rf}$ est compatible avec $W(I')_{\rf}$, $W(I)_{\rf}$, il s'ensuit que la flèche horizontale $\alpha$
\msp{page 71}
est compatible avec les relations d'équivalence définies par
\[
\Hom(F',G') \longrightarrow \Hom_{W(I')}(F',G') \quad \text{resp.} \ \ \Hom(u_*F',u_*G') \longrightarrow \Hom_{W(I)}(u_*F',u_*G')\ .
\]
Et il en est de même de la flèche $\mu : f \longmapsto f \circ i$, trivialement. Donc la flèche d'adjonction $\Hom(u^*F,G') \longrightarrow \Hom(F,u_*G')$ est compatible avec les relations $R$, $R'$, ce qui achève de prouver que $(**)$, page 68, induit une bijection $(*)$, page 68.

Ceci montre donc

\begin{enumerate}
\item[1°)] que pour un objet fixé $G'$ dans $\HO_W(I')$ (qu'on peut toujours supposer, à isomorphisme près, provenir de $\M(I')_{\rf}$), le foncteur
\[
F \ \longmapsto\ \Hom_{\HO_{W}(I')}(u^*_W F,G')
\]
est représentable -- \textit{i.e.} le foncteur $u^*_W$ admet un adjoint à gauche, et
\item[2°)] que cet adjoint à gauche se calcule comme $u_*(G')$ (quand l'argument $G'$ est pris fibrant) -- vu que l'homomorphisme $({*}{*}{*})$, page 69, est \textit{fonctoriel en $F$}. (A posteriori il s'ensuit qu'il est aussi fonctoriel en $G'$ \dots)
\end{enumerate}

\printedpage{40}
% head: A. Grothendieck : Dérivateurs, Ch. XIII, Édition des Universités de Montpellier II et Paris VII
% foot: Version 25 janvier 2004

\msp{page 72}

\dgrblock{dgrb721}{Pointwise right images require fibrant restriction}
\begin{dcomment}{exposition}{p71}{Pointwise right images require fibrant restriction}
Grothendieck now wants more than an abstract right adjoint: he wants to compute its value at an object by a homotopy limit over a comma category. The comparison is canonical, but it becomes an isomorphism only after the same replacement is known to be suitable on both sides. This is the content of the invocation of local isomorphisms on p.~72.

Fix $s\in I$, let $q:I'/s\to I'$ be the projection, and write $p_s:I'/s\to e$. If $X'\to RX'$ is fibrant replacement and $q^*$ preserves fibrant objects, then $q^*RX'$ is a fibrant replacement of $q^*X'$, because restriction always preserves the objectwise weak equivalences. The ordinary pointwise formula therefore identifies
\[
 \bigl(\mathbf Ru_*X'\bigr)_s
   =\bigl(u_*RX'\bigr)_s
   \cong\varprojlim_{I'/s}q^*RX'
   \cong\mathbf R(p_s)_*q^*X'.
\]
The last step uses admissibility of $p_s$. Thus the comparison is computed on one common representative, rather than by comparing two unrelated fibrant replacements.

For clarity, there is a simple sufficient condition for the fibrancy assertion. The functor $q$ is a discrete fibration: an arrow $j\to j'$ in $I'$ and an object $(j',u(j')\to s)$ have a unique lift with source $(j,u(j)\to u(j')\to s)$. In our contravariant convention, when the necessary coproducts exist, its left Kan extension has the formula
\[
                           (q_!A)_j\cong\coprod_{x\in q^{-1}(j)}A_x.
\]
Coproducts of trivial cofibrations have the left lifting property against every fibration, so $q_!$ preserves the objectwise trivial cofibrations. Its right adjoint $q^*$ consequently preserves fibrations and fibrant objects. This is the contravariant version of the discrete-opfibration argument in \citet[Lemma 1.31]{Groth2013}.

If $I$ is a poset, each fiber has at most one element, so only the initial object is needed for this coproduct formula. In the large cocomplete model examples, all small fibers are allowed. In a setting with only finite limits and finite colimits, an arbitrary infinite fiber requires an additional existence hypothesis; it must not be supplied by notation. This is why the general admissibility result and the pointwise formula have been stated separately in the commentary. The finite-poset corollary satisfies the needed conditions, and the usual combinatorial model setting satisfies them for every small functor.

Modern derivator theory packages this calculation as the pointwise Kan-extension axiom. The square of shapes contains the canonical natural transformation $uq\Rightarrow s p_s$; its mate is exactly the displayed comparison. Thinking in terms of that mate explains its compatibility with composition and base change, while the fibrant-representative calculation proves that it is invertible in the present models.

The same proof gives a useful class of base-change squares, with a condition that can be checked directly. Let $v:I_1\to I$ be a discrete fibration, form the strict pullback $I'_1=I'\times_I I_1$, and denote its projections by $v':I'_1\to I'$ and $u_1:I'_1\to I_1$. Assume that the required right Kan extensions and the coproducts over the fibers exist. For each $t\in I_1$ the natural functor
\[
                         I'_1/t\longrightarrow I'/v(t)
\]
is an isomorphism. Indeed, an object of the target is an arrow $u(i')\to v(t)$. Its unique lift along $v$ is an arrow $t'\to t$ with $v(t')=u(i')$, and therefore supplies the object $((i',t'),t'\to t)$ of the source. The uniqueness of lifts gives the inverse also on morphisms. This proves the ordinary base-change identity by the pointwise limit formula.

The pullback $v'$ is again a discrete fibration, so the preceding coproduct argument shows that $(v')^*$ preserves fibrant objects. Evaluate both sides of the ordinary identity on a fibrant replacement of $X'$. They then compute the two derived functors, proving the canonical isomorphism
\[
                      v^*\mathbf Ru_*X'
                     \cong\mathbf R(u_1)_*(v')^*X'.
\]
This is the same discrete-fibration mechanism used in the comma calculation, applied to a family of objects at once; it is part of the standard pointwise derivator calculus cited above.

A strict pullback square without this hypothesis need not have such an isomorphism. Take $I=(0\to1)$, let $I'=\{0\}$ and $I_1=\{1\}$ be the two endpoint inclusions, and use complexes over a field. Their pullback is empty. For $X'=k[0]$, the left side evaluated at $1$ is $k[0]$, because $\{0\}/1$ is terminal. The right side is zero, the limit of the empty diagram in complexes. Thus the base-change map cannot be an isomorphism. The pointwise axiom guarantees the designated comma squares; additional squares require the corresponding exactness argument.
\end{dcomment}


Le dernier énoncé dans 2° du théorème-scholie résulte du calcul de $u_*^W$, compte tenu que $I'/s \longrightarrow I'$ est un isomorphisme local, donc le foncteur image inverse envoie
\[
\M(I')_{\rf} \longrightarrow \M(I'/s)_{\rf},
\]
donc si $G'$ est dans $\M(I')_{\rf}$,
\[
u_*(G')(s) \ \simeq\ p_{s\,*}(G'|I'/s)
\]
est un calcul de $p^W_{s\,*}(G'|I'/s)$.

Cela achève la démonstration du théorème-scholie sur les images directes.

\textbf{Commentaires.}

1) Ce théorème me paraît loin d'être satisfaisant, c'est un pis-aller. J'avais espéré, avec les hypothèses du paragraphe 3, qu'on aurait l'existence des $u_*^W$ pour \textit{toute} flèche $u : I' \longrightarrow I$ dans Cat. Je doute à présent que ce soit le cas, mais n'ai pas encore réussi de faire un contre-exemple. Les cas test les plus importants pour la question d'admissibilité d'une catégorie $I$ me paraissent les suivants :

\begin{enumerate}
\item[\textit{a})] $I$ catégorie discrète infinie (disons dénombrable). La question revient ici à celle-ci : soient $f_i : X_i \longrightarrow Y_i$ dans $W$, avec $X_i$, $Y_i$ fibrants, \textit{i.e.} dans $F$ \texttt{[\textit{i.e.}} $X_i \longrightarrow e$\texttt{,} $Y_i \longrightarrow e$ \texttt{dans} $F$\texttt{]}. Alors $f = \prod f_i : \prod X_i \longrightarrow \prod Y_i$ est-il dans $W$ ?

\msp{page 73}

\dgrblock{dgrb722}{Products and towers have a positive answer}
\begin{dcomment}{status}{p72}{Products and towers have a positive answer}
Grothendieck's test cases isolate two ways in which an infinite limit could lose a weak equivalence. The first takes infinitely many independent factors; the second imposes infinitely many compatibility equations. Under the closed model hypotheses he explicitly mentions, both questions have affirmative answers. Their difficulty is the availability of suitable replacements, not a failure of weak equivalence to survive the limit of a fibrant diagram.

For a set $S$, equip $\mathcal M^S$ with the product structure. The diagonal functor preserves cofibrations and trivial cofibrations componentwise. Its right adjoint, whenever it exists, is $\prod_S$. The Brown argument therefore shows that a product of weak equivalences between fibrant objects is a weak equivalence. This remains true for uncountable $S$ whenever the required products exist. No assumption about exact filtered colimits is used. The derived product is consequently computed by the product of fibrant replacements; the same conclusion appears explicitly in \citet[Example 1.3.11]{Hovey1999}.

For the tower shape, use the matching-object description already seen in finite posets. A tower
\[
                         X_0\longleftarrow X_1\longleftarrow X_2\longleftarrow\cdots
\]
is fibrant precisely when $X_0$ is fibrant and every bonding map is a fibration; all other $X_n$ are then fibrant by composition. The recursive lifting proof has one step for each natural number and defines a compatible lift on the whole tower. Taking the limit is right adjoint to the constant-tower functor, which preserves the objectwise cofibrations and trivial cofibrations. Brown's lemma again gives invariance under levelwise weak equivalences between these towers. In the ordinary complete model setting this is also the inverse Reedy construction.

The manuscript correctly observes that the tower assertion implies the countable product assertion. For a countable family $(X_n)$ of fibrant objects, put $P_r=\prod_{0\leq n\leq r}X_n$. The projection $P_{r+1}\to P_r$ is a pullback of $X_{r+1}\to e$, hence a fibration. A family of weak equivalences induces weak equivalences between the finite products, by the finite case; the inverse limit of the resulting map of towers is the original map of infinite products. This reduction proves the stated implication, although the direct product adjunction is shorter once the general lemma is recognized.

The affirmative answer does not mean that ordinary inverse limits preserve all levelwise weak equivalences. The following example shows exactly where the bonding-map condition matters. It is also a useful safeguard when replacing an abstract existence theorem by an explicit tower of complexes.
\end{dcomment}


\item[\textit{b})] I catégorie % [sic] I (upright roman) -> $I$, italic everywhere else for the category
\[
0 \longrightarrow 1 \longrightarrow 2 \longrightarrow 3 \longrightarrow \cdots,
\]
\textit{i.e.} la catégorie ordonnée associée à l'ensemble ordonné $\mathbf{N}$, pour la relation d'ordre ordinaire $m \leq n$. Donc $\M(I)$ est formée des systèmes \textit{projectifs} dans $\M$, indexés par $\mathbf{N}$. On prouve facilement (en procédant comme dans le corollaire du lemme p. 63) que $\M(I)_{\rf}$ est formée des systèmes projectifs
\[
X_0 \longleftarrow X_1 \longleftarrow X_2 \longleftarrow \cdots \longleftarrow X_n \xleftarrow{\ p_n\ } X_{n+1} \longleftarrow \cdots
\]
tels que les $X_i$ soient dans $\M_{\rf}$, et les $p_n$ dans $F$. La question est donc : si $X_\bullet$, $Y_\bullet$ sont fibrants et $f_\bullet : X_\bullet \longrightarrow Y_\bullet$ est telle que $f_i \in W$ pour $i \in \mathbf{N}$, alors a-t-on $\varprojlim f_i : X_\infty \longrightarrow Y_\infty$ dans $W$ ?

Si la reponse est affirmative pour $\mathbf{N}$ (pour un $W$ et une quillénisation donnée), elle % [sic] reponse -> réponse, accent missing
l'est aussi pour toute catégorie discrète dénombrable, comme on le voit aussitôt. Mais je ne vois pas comment déduire même \textit{a}) des axiomes, disons pour un triplet de Quillen clos (closed model category) tel que $\M$ soit de plus stable par $\varinjlim$ filtrantes, que celles-ci soient exactes, et qu'il y ait
\printedpage{41}
% head: A. Grothendieck : Dérivateurs, Ch. XIII, Édition des Universités de Montpellier II et Paris VII
% foot: Version 25 janvier 2004
\msp{page 74}
un petit ensemble de générateurs dans $\M$ (donc $\M$ est accessible), que $W$ soit stable par $\varinjlim$ filtrantes. Je devrais demander à \textsc{Thomason} s'il connait un contre-exemple. % [sic] connait -> connaît, circumflex missing

\dgrblock{dgrb723}{An inverse limit of acyclic complexes can have cohomology}
\begin{dcomment}{examples}{p73}{An inverse limit of acyclic complexes can have cohomology}
The insistence on fibrant towers is mathematically necessary, even in the most familiar algebraic model. Fix a prime $p$ and, for $n\geq1$, form the cochain complex
\[
 A_n=\bigl[\mathbb Z\xrightarrow{p^n}\mathbb Z
                   \xrightarrow{\bmod p^n}\mathbb Z/p^n\bigr]
\]
in degrees $0,1,2$. It is acyclic: the first map is injective, the image of that map is the kernel of the second, and the second is surjective. Define $A_{n+1}\to A_n$ by multiplication by $p$ in degree $0$, the identity in degree $1$, and reduction modulo $p^n$ in degree $2$. These maps commute with both differentials. The maps of towers $A_n\to0$ are therefore quasi-isomorphisms at every level.

Their ordinary inverse limit is readily computed. A compatible sequence in degree $0$ satisfies $a_n=p a_{n+1}$. Its first coordinate is divisible by every power of $p$, and hence vanishes; then every coordinate vanishes. The limit in degree $1$ is $\mathbb Z$, and the limit in degree $2$ is $\mathbb Z_p$. Thus
\[
               \varprojlim_n A_n=\bigl[0\longrightarrow\mathbb Z
                                          \longrightarrow\mathbb Z_p\bigr].
\]
The map $\mathbb Z\to\mathbb Z_p$ is injective but not surjective. For example, $\mathbb Z_p$ is uncountable, since a sequence of digits in $\{0,\ldots,p-1\}$ uniquely determines a $p$-adic integer, while $\mathbb Z$ is countable. Consequently
\[
              H^2\!\left(\varprojlim_n A_n\right)=\mathbb Z_p/\mathbb Z\ne0.
\]
The inverse limit of the levelwise quasi-isomorphism $A_\bullet\to0$ is not a quasi-isomorphism.

There is no conflict with the positive result on the preceding page. In the projective model structure on unbounded complexes of abelian groups every object is fibrant, but a fibration is degreewise surjective. The bonding map of this tower is multiplication by $p$ in degree $0$, so it is not a fibration. Individual fibrancy has not made the \emph{diagram} fibrant. A fibrant tower replacement would have an acyclic inverse limit by the Brown argument.

This example separates three statements which can otherwise look deceptively similar: ordinary inverse limit is not exact; inverse limit need not preserve levelwise quasi-isomorphisms of arbitrary towers; inverse limit does preserve weak equivalences between fibrant towers in the model structure under discussion. Only the third statement is the admissibility problem posed in the manuscript. Its positive answer is compatible with the failure exhibited here.
\end{dcomment}


\item[\textit{c})] Les deux exemples précédents étaient des catégories ordonnées. Comme exemple typique de catégorie non-ordonnée, où je ne vois pas ce qui se passe, je signale $I = B_G$, $G$ un groupe (donc un seul objet, avec un groupe d'automorphismes $G$). Aussi le cas de $\boldsymbol{\Delta}[1] = \{\boldsymbol{\Delta}_0, \boldsymbol{\Delta}_1\}$ (sous-catégorie pleine de $\boldsymbol{\Delta}$) n'est pas clair pour moi.
\end{enumerate}

\dgrblock{dgrb724}{A group action detects the difference between fixed and homotopy fixed points}
\begin{dcomment}{examples}{p74:Les deux exemples précédents étaient des catégories ordonnées.}{A group action detects the difference between fixed and homotopy fixed points}
Grothendieck introduces $B_G$ because automorphisms are absent from the finite-poset argument. The limit of a $G$-diagram is a fixed-point object. Its derived version measures the coherence of invariance under the whole action, and a single fibrant value does not guarantee that the action diagram is injectively fibrant. The algebraic example below makes the difference explicit and at the same time answers the boundedness question from \S5.

Let $G=C_p=\langle g\mid g^p=1\rangle$, let $k$ be a field of characteristic $p$, and give $k$ the trivial action. Using inversion to change a right action into a left action if necessary, diagrams of complexes on $B_G$ are complexes of $R=k[G]$-modules. With $t=g-1$, one has $R\cong k[t]/(t^p)$. The trivial module has the free resolution
\[
 \cdots\xrightarrow{t}R\xrightarrow{t^{p-1}}R
       \xrightarrow{t}R\longrightarrow k\longrightarrow0.
\]
Its exactness follows immediately from
\[
       \ker(t:R\to R)=(t^{p-1}),\qquad
       \ker(t^{p-1}:R\to R)=(t).
\]
Applying $\operatorname{Hom}_R(-,k)$ makes every differential zero, since $t$ acts by zero on $k$. Tensoring with $k$ does the same. Therefore
\[
 \operatorname{Ext}^n_R(k,k)\cong k,
 \qquad \operatorname{Tor}^R_n(k,k)\cong k
 \qquad(n\geq0).
\]
For the projection $p:B_G\to e$ this gives
\[
 \mathbf Rp_*k=\mathbf R\operatorname{Hom}_{k[G]}(k,k),
 \qquad \mathbf Lp_!k=k\otimes^{\mathbf L}_{k[G]}k.
\]
The first object has nonzero cohomology in every nonnegative degree, and the second in every nonpositive degree. Both exist in $D(k)$. Neither belongs to $D^{\mathrm b}(k)$, although the input is concentrated in degree zero. Finite cohomological dimension is therefore a boundedness condition in this example, not a prerequisite for the unbounded operations.

Ordinary invariants of the original diagram give only $k[0]$. Moreover, over a field its underlying complex is already fibrant. The additional cohomology arises because injective fibrancy in the diagram category is a condition on the $k[G]$-complex, not merely on the underlying $k$-complex. Replacing that diagram by a $K$-injective $k[G]$-complex and taking ordinary invariants computes $\mathbf Rp_*$. The Brown lemma proves invariance between such fibrant diagrams, so the presence of a nontrivial automorphism group creates no exception to admissibility.

Exactly the same general argument applies to the full subcategory of $\boldsymbol\Delta$ on $[0]$ and $[1]$ mentioned next in the manuscript. It is not the walking arrow: it contains all maps between those two ordinals, including the two sections and the retraction. It has two objects and seven morphisms: the four Hom sets have respectively $1,2,1,3$ elements. Thus its ordinary limit uses only the finite limits of $\mathbf M0$. Its relations matter when constructing a coherent diagram, but the assumed uniform structure and the lifting-adjunction proof already give admissibility; no further infinite-completeness assumption is needed for this test.
\end{dcomment}


2) On aimerait pouvoir traiter de façon duale l'existence des $u_!^W$, à défaut de mieux. Dans l'hypothèse de départ du paragraphe 4 sur $(\M,W)$, cela amène donc à changer de tactique. Cette fois, on devra supposer l'existence d'une quillénisation close \textit{exacte à droite} pour $W$ (ce qui est plus exigeant, semble-t-il, qu'une telle structure exacte à gauche, dans le cas envisagé), donc
\[
TF \ = \ F \cap C,
\]
\msp{page 75}
et de plus encore $C$, $TC$ tels qu'on ait
\[
C \ = \ \widetilde{C_0}, \qquad TC \ = \ \widetilde{TC_0}
\]
avec $C_0$, $TC_0$ petits (\footnote[38]{On dit alors que $C$, $TC$ (en tant que parties closes à droite \texttt{[plutôt à gauche]}) sont “accessibles”.}). On posera maintenant
\[
\begin{aligned}
F(I) \ &= \ \{f \in \Fl \M(I) \ \mid \ f_s \in F \quad\ \ \forall\, s \in \Ob I\}\\
TF(I) \ &= \ \{f \in \Fl \M(I) \ \mid \ f_s \in TF \ \ \forall\, s \in \Ob I\},
\end{aligned}
\]
puis
\[
\begin{aligned}
C(I) \ &= \ TF(I)^*\\
TC(I) \ &= \ F(I)^*\ .
\end{aligned}
\]
L'espoir (que je devrai vérifier ultérieurement), c'est qu'on obtient ainsi une quillénisation uniforme à droite de $W$, sur Cat tout entier. Le seul problème, c'est la validité des axiomes de factorisation de $f \in \Fl \M(I)$ en $f = pi$, avec $i \in TC$, $f \in F$, et aussi avec $i \in C$, $f \in TF$. % [sic] f \in F, f \in TF -> p \in F, p \in TF, the factorisation is f = pi
J'ai bon espoir que c'est vrai -- cela se réduit essentiellement à voir que $(C(I),TF(I))$ et $(TC(I),F(I))$, soit $(\Phi(I),\Psi(I))$, satisfont en plus de $\Phi(I) = \Psi(I)^*$, à $\Psi(I) = \Phi(I)_*$ (\textit{i.e.} $\Psi(I)$ est clos à droite), et que $\Phi(I)$ est accessible (tout comme $\Phi = \Phi(e)$). Si ça marche, on aura donc l'existence
\printedpage{42}
% head: A. Grothendieck : Dérivateurs, Ch. XIII, Édition des Universités de Montpellier II et Paris VII
% foot: Version 25 janvier 2004
\msp{page 76}
de $u_!$ pour $u$ un foncteur entre catégories ordonnées finies, et sous des conditions légèrement plus générales. (P. ex. $u : I' \longrightarrow I$ avec $I$ ordonnée, $I'$ ordonnée et à espace sous-jacent noethérien et sobre.) Peut-être que cette fois, l'hypothèse que $W$ est stable par $\varinjlim$ filtrantes va quand même servir! Ainsi, le cas des catégories discrètes, et des catégories filtrantes à gauche comme $\mathbf{N}^\circ$, ne va plus poser de problèmes à cause de cela. Il y a donc de l'espoir! Tout ensemble ordonné qui est $\varinjlim$ filtrante de sous-ensembles ordonnés \textit{fermés} noethériens et sobres, est alors $W$-admissible à droite. Mais qu'en est-il de catégories de type $\Gamma = B_G$ p. ex. ? ? ?

\dgrblock{dgrb725}{The projective construction supplies the dual theory}
\begin{dcomment}{later}{p75}{The projective construction supplies the dual theory}
Grothendieck's change of tactic is the right one for left images. To derive a colimit, fibrations should be defined objectwise, so restriction becomes a right Quillen functor and its left adjoint acts on cofibrant diagrams. The extra small generating sets are needed to construct these projective diagram structures. Once they have been constructed, the dual Brown argument supplies homotopy invariance for every existing left Kan extension, rather than only for finite posets.

Here is the construction with its size hypothesis visible. Suppose that $\mathcal M$ is a locally presentable, cofibrantly generated model category, with generating cofibrations $C_0$ and generating trivial cofibrations $T_0$. For $i\in I$, evaluation has a left adjoint $L_i$, given in the contravariant convention by
\[
                         (L_iA)_j=\coprod_{I(j,i)}A.
\]
Use the sets $\{L_ic:i\in I,c\in C_0\}$ and $\{L_it:i\in I,t\in T_0\}$. By the evaluation adjunction, their right lifting classes are exactly the objectwise trivial fibrations and the objectwise fibrations, respectively. The sets really are small: $I$ and the two generating collections are small. For a sufficiently large regular cardinal the domains are presentable, because evaluation preserves filtered colimits and
\[
 \operatorname{Hom}_{\mathcal M(I)}(L_iA,X)
                        \cong\operatorname{Hom}_{\mathcal M}(A,X_i).
\]
The small object argument therefore gives both desired factorizations.

One also checks the weak equivalences in the transferred structure. A relative cell complex built from the $L_it$ is componentwise a composite of pushouts of coproducts of trivial cofibrations. The base trivial cofibrations form a left lifting class and are closed under these operations, so the resulting map is an objectwise trivial cofibration and hence an objectwise weak equivalence. Taking retracts gives the same conclusion for every map in $TC(I)$. Moreover $TC(I)\subset C(I)$, because the objectwise trivial fibrations are objectwise fibrations. If $f\in C(I)\cap W(I)$, factor it as $f=pj$ with $j\in TC(I)$ and $p\in F(I)$. Two-out-of-three gives $p\in W(I)$, so $p\in TF(I)$ by the objectwise definitions. The lifting property of $f$ against $p$ makes $f$ a retract of $j$, hence $f\in TC(I)$. This identifies the required acyclic class and completes the model-structure verification. It is the projective transfer argument sought in the paragraph following the two formulas on p.~75. See \citet[Theorem 1.28 and Lemma 1.29]{Groth2013}; the small object argument itself is recorded in \citet[Theorem 2.1.14]{Hovey1999}.

For $u:I'\to I$, restriction now preserves the objectwise fibrations and trivial fibrations. Adjointness makes $u_!$ preserve cofibrations and trivial cofibrations. The dual Brown lemma shows that it preserves weak equivalences between cofibrant diagrams, and the derived adjunction is
\[
                              \mathbf Lu_!\dashv u^*.
\]
The hypotheses of local presentability and cofibrant generation make all small diagram shapes available, including $B_G$. An objectwise construction of fibrations does not mean that cofibrant diagrams are objectwise cofibrant with no further condition; the latter would miss precisely the relations that the free diagrams $L_iA$ impose.

Two slips in the source obscure this otherwise correct programme. Exactness on the right is $TF=F\cap W$, not $F\cap C$. In the factorizations $f=pi$, the membership assertions must refer to $p\in F$ or $p\in TF$, not to $f$. They are retained as printed and explained in the correction register. The role of filtered colimits also deserves precision: in the projective transfer just described, closure of trivial cofibrations under transfinite composites follows from their lifting property. A separate assertion that all weak equivalences are closed under filtered colimits is not what proves that step.
\end{dcomment}


\msp{page 77}

\dgrblock{dgrb726}{Why the homotopy problem returns to SGA 4}
\begin{dcomment}{context}{p76}{Why the homotopy problem returns to SGA 4}
Grothendieck turns to cardinal filtrations because a construction on every diagram category needs a controlled supply of small pieces. The issue is already visible in the proposed right uniform quillenisation: a lifting condition must be generated by a set of maps, and the factorisations built from that set must remain available after changing the indexing category. The return to SGA~4, I, \S9 is therefore part of the homotopical argument. It supplies a way of bringing a large category within reach of set-indexed constructions.

The later letter to Thomason makes this connection explicit. In its opening pages Grothendieck describes his renewed work on accessible categories and accessible parts as a rebuilding of the provisional account in SGA~4. In the next section of the letter, he distinguishes the relative category $(\mathcal M,W)$, which specifies the homotopy theory, from the cofibrations and fibrations used to construct its operations \citep[pp.~2--4 of the electronic edition]{Grothendieck1991Letter}. The present chapter places these two concerns next to one another: first find enough replacements to derive direct images; then explain the smallness that makes such replacements possible. This is also why the size of a \emph{diagram}, rather than just the sizes of its values, becomes decisive near the end of \S7.

There are two historical cautions. The word ``accessible'' here follows the cardinal conventions of SGA~4: ``grand devant $\pi$'' bounds every subset of size at most $\pi$, and thus corresponds to $\pi^+$-directedness in the regular-cardinal terminology now customary. Also, the closing definition tests membership in a part against \emph{every} perfectly adapted presentation. It is stronger than merely exhibiting the part as a category of sufficiently directed colimits of small objects. The distinction is mathematically effective: the counterexamples below separate those conditions even in a powerset ordered by inclusion. The later development of accessibility does not license us to identify formulations whose quantifiers differ.
\end{dcomment}


\dgrblock{dgrb727}{What has been settled before the cardinal detour}
\begin{dcomment}{status}{p76}{What has been settled before the cardinal detour}
The final questions of \S6 can now be separated by the task they perform. For the uniform closed exact structures assumed in the theorem-scholie, the Brown lemma settles admissibility of every functor with the requisite ordinary right adjoint; its dual settles left admissibility. Products, towers, group actions and the truncated simplex category are all included. The finite-poset proof remains a direct construction of the finite calculus, while the noetherian sober Alexandrov condition adds no infinite posets.

The next question is existence of those uniform structures. In the usual combinatorial setting, injective and projective diagram model structures, derived adjunctions, and the pointwise comma-category formulas furnish a strong derivator on all small categories. This is the later theorem corresponding most directly to Grothendieck's programme. In smaller categories such as bounded complexes, one retains finite operations and must check each infinite ordinary limit or colimit separately. The bounded-product example above prevents a blanket completeness assertion.

Section 7 addresses the remaining constructive burden: how a large lifting class or a large collection of objects can be controlled by a small family at a sufficiently large cardinal. That question is substantial even though the preceding homotopy-invariance doubts have a formal answer. It is also where the older SGA~4 language of accessibility and the later language of presentability need to be compared with particular care. The change of subject is therefore motivated by the production of factorizations, rather than by a need to repair Brown's argument.
\end{dcomment}


\section*{7\quad Filtrations cardinales et parties accessibles d'une grosse catégorie}

\dgrblock{dgrb728}{A size filtration and a presentability rank}
\begin{dcomment}{exposition}{p77:Filtrations cardinales et parties accessibles d'une grosse catégorie}{A size filtration and a presentability rank}
Grothendieck returns to cardinal estimates because the constructions of model structures need a controlled supply of small objects. A replacement construction may be performed object by object and still fail to be controlled uniformly over a category. The filtration is intended to provide that control: a large object is assembled from bounded pieces, while bounded constructions on bounded pieces remain bounded. The categorical purpose of the estimates should remain visible through the set theory.

There is a one-cardinal shift in the language. Throughout these pages, ``grand devant $\pi$'' means that every subset of cardinal at most $\pi$ has an upper bound. In modern terminology such a nonempty poset is $\pi^+$-directed, where $\pi^+$ is the successor cardinal. Write $\kappa=\pi^+$ when translating the assertions. An object called $\pi$-accessible in the sense of SGA~4 is correspondingly $\kappa$-presentable: its covariant Hom functor preserves these directed colimits. This meaning of ``accessible object'' must also be kept distinct from the new meaning of ``accessible part'' introduced on manuscript page~86. The latter concerns a full subcategory and its position inside the ambient category.

The three filtration conditions do not say that every object of $\Filt^\pi(E)$ is $\pi^+$-presentable. They say that this stage is essentially small, is closed under the specified bounded colimits, and generates $E$ under sufficiently directed colimits. Preservation of such colimits by Hom is an additional hypothesis. This explains why the hypothesis $\Filt^\pi(E)\subset E_\pi$, initially left tacit, becomes necessary on pages~90--94. Under suitable accessibility hypotheses, SGA~4 supplies useful cardinals at which size and presentability do coincide; see \citet[I, 9.16--9.18]{SGA4I1972}. It does not identify the two notions at every index of an arbitrary filtration.

The universe convention is also doing real work. ``Essentially small'' means that a skeleton belongs to the fixed universe; it does not mean that there are finitely or countably many objects. Once a skeleton is chosen, the relevant comma categories are small because $E$ is a $\mathfrak U$-category. All diagrams below remain small relative to this universe. No assertion that the whole large category has a set of objects in that same universe is needed.

Finally, the filtration is exhaustive even though exhaustiveness is not listed separately. Choose the presentation from condition~(c) at $\pi_0$, and choose $\rho\geq\pi_0$ larger than the cardinality of its indexing poset. Its terms lie in $\Filt^\rho(E)$, and condition~(b) places its colimit in the same stage. Thus every object eventually enters the filtration. This short argument is the content of footnote~39.
\end{dcomment}


Je me réfère à la théorie des foncteurs et catégories accessibles de SGA 4, I, §9 (p. 138-178). Je me donne une grosse catégorie (toujours une $\mathfrak{U}$-catégorie) $E$ avec une \textit{filtration cardinale} (croissante)
\[
\Bigl(\Filt^{\pi}(E)\Bigr)_{\pi \geq \pi_0}, \qquad \pi_0 \text{ un cardinal infini fixé,}
\]
par les sous-catégories strictement pleines $\Filt^{\pi}(E)$. On a \texttt{[SGA 4, I, 9.12]}

\begin{enumerate}
\item[\textit{a})] $\Filt^{\pi}(E)$ essentiellement petite, \textit{i.e.} équivalente à une petite catégorie (\textit{i.e.} à une catégorie appartenant à $\mathfrak{U}$).
\item[\textit{b})] $E$ satisfait $L_{\pi_0}$ (\textit{i.e.} $E$ est stable par $\varinjlim$ de type $I$, quand $I$ est un ensemble ordonné grand devant $\pi_0$), et $\Filt^{\pi}(E)$ est stable par les limites inductives filtrantes indexées par $I$ grand devant $\pi_0$, et tel que $\card I \leq \pi$.
\item[\textit{c})] Pour tout $\pi \geq \pi_0$, et pour tout $X \in \Ob E$, on a
\[
X \ \simeq\ \varinjlim_{i \in I} X_i,
\]
où les $X_i$ \texttt{[sont]} dans $\Filt^{\pi}(E)$, et où $I$ est un ensemble ordonné grand devant $\pi$. De plus, si $X \in \Ob \Filt^{\pi'}(E)$, avec $\pi' > \pi$, on peut prendre $I$ tel que $\card I \leq \pi'^{\pi}$ (donc si $\pi' = 2^{c'}$ avec $c' \geq \pi$, d'où $(\pi')^{\pi} = (2^{c'})^{\pi} = 2^{c'\pi} = 2^{c'} = \pi'$, on aura $\card I \leq \pi'$). (\textbf{NB} Si on admet l'hypothèse du continu généralisée, on aura toujours $\pi'^{\pi} = \pi'$, donc $\card I \leq \pi'$.) (\footnote[39]{NB \textit{b}) et \textit{c}) impliquent que $E = \bigcup_{\pi \geq \pi_0} \Filt^{\pi}(E)$.})
\end{enumerate}

\dgrblock{dgrb729}{The useful exponent bound and the singular-cardinal exception}
\begin{dcomment}{examples}{p77:Si on admet l'hypothèse du continu généralisée}{The useful exponent bound and the singular-cardinal exception}
The bound $(\pi')^\pi$ measures how many small pieces may be needed to describe an object of size $\pi'$. It should be retained until an actual cardinal identity permits its replacement by $\pi'$. The calculation for a power-set cardinal is correct: if $\pi'=2^{c'}$ and $c'\geq\pi$ are infinite, then
\[
 (\pi')^\pi=(2^{c'})^\pi=2^{c'\cdot\pi}=2^{c'}=\pi'.
\]
Thus arbitrarily large useful choices are available without a continuum hypothesis.

The subsequent claim about the generalized continuum hypothesis is false as written. If $\mu$ is singular and $\operatorname{cf}(\mu)\leq\pi<\mu$, then $\mu^\pi>\mu$, independently of GCH. Here is the elementary obstruction. Choose increasing infinite cardinals $\mu_\alpha<\mu$, for $\alpha<\delta=\operatorname{cf}(\mu)$, with supremum $\mu$. K\"onig's diagonal argument gives
\[
 \mu=\sum_{\alpha<\delta}\mu_\alpha
       <\prod_{\alpha<\delta}\mu_\alpha^+
       \leq\mu^\delta\leq\mu^\pi.
\]
For completeness, the strict inequality follows by partitioning a purported list of all elements of the product into pieces $S_\alpha$ of size $\mu_\alpha$. In coordinate $\alpha$, choose a value outside the $\alpha$-coordinates of $S_\alpha$; the resulting element is absent from every piece of the list. Taking $\mu=\aleph_\omega$ and $\pi=\aleph_0$ already disproves the printed assertion.

Under GCH the conclusion does hold when $\pi<\operatorname{cf}(\pi')$. In that case each map $\pi\to\pi'$ has bounded image. For every infinite $\nu<\pi'$, GCH gives
\[
 \nu^\pi\leq(2^\nu)^\pi=2^{\max(\nu,\pi)}
 =\max(\nu,\pi)^+\leq\pi',
\]
and the union of $\pi'$ such sets has size at most $\pi'$. The constant functions give the reverse inequality. The missing cofinality condition matters precisely at singular cardinals. Nothing later in the construction requires the erroneous unconditional simplification: the bound $(\pi')^\pi$, and the unconditional power-set choices, suffice.
\end{dcomment}


\printedpage{43}
% head: A. Grothendieck : Dérivateurs, Ch. XIII, Édition des Universités de Montpellier II et Paris VII
% foot: Version 25 janvier 2004

\msp{page 78 barrée}

\msp{page 79}

Considérons la représentation de \textit{c}),
\[
X \ = \ \varinjlim_{I} X_i \qquad\qquad
\left\{
\begin{aligned}
& X_i \in \Filt^{\pi},\\
& I \text{ grand devant } \pi.
\end{aligned}
\right.
\]
Soit
\[
\tilde{I} \ = \ \{I' \subset I \mid \card I' \leq \pi,\ \ I' \text{ grand devant } \pi_0\},
\]
muni de la relation d'ordre d'inclusion. Pour
\[
I' \ \in\ \tilde{I}, \qquad \text{soit}\ \ X_{I'} = \varinjlim_{I'} X_i\ ,
\]
cette limite existe par $L_{\pi_0}$ (hypothèse \textit{b}), elle est dans $\Filt^{\pi}$ par \textit{b}), donc on a un système inductif
\[
(X_{I'})_{I' \in \tilde{I}}, \qquad X_{I'} \in \Filt^{\pi},
\]
avec
\[
\card \tilde{I} \ \leq\ (\card I)^{\pi},
\]
et si $X \in \Filt^{\pi'}$ $(\pi' > \pi)$, prenant
\[
\card I \ \leq\ \pi'^{\pi},
\]
on aura $\card \tilde{I} \leq (\pi'^{\pi})^{\pi} = \pi'^{\pi^2} = \pi'^{\pi}$, donc
\[
\card \tilde{I} \ \leq\ \pi'^{\pi}.
\]
Je dis que $\tilde{I}$ est lui aussi \textit{grand devant $\pi$,} tout comme $I$. En effet,
\msp{page 80}
soit
\[
(I_\alpha)_{\alpha \in A}, \qquad I_\alpha \in \tilde{I}, \qquad \card A \ \leq\ \pi,
\]
prouvons qu'il existe $I' \in \tilde{I}$ qui majore les $I_\alpha$, \textit{i.e.} qui contient $I'_0 = \bigcup_{\alpha \in A} I_\alpha$. On aura $\card I'_0 \leq \pi^2 = \pi$ (car $\card A \leq \pi$ et $\card I_\alpha \leq \pi$ pour tout $\alpha$), et il suffit de voir ceci :

\texttt{[En fait, l'assertion est évidente : comme} $\card I'_0 \leq \pi$\texttt{, et comme} $I$ \texttt{est grand devant} $\pi$\texttt{,} $I'_0$ \texttt{admet un majorant} $i'_0$ \texttt{dans} $I$\texttt{, et on peut prendre} $I' = I'_0 \cup \{i'_0\}$\texttt{.]}

\textbf{Lemme 1 :} \textit{Soient $\pi_0$, $\pi$ deux cardinaux, avec $\pi \geq \pi_0$, $I$ un ensemble ordonné grand devant $\pi$, et $I'_0$ une partie de $I$ de cardinal $\leq \pi$. Alors il existe une partie $I'$ de $I$ satisfaisant les conditions :} (\footnote[40]{Pourvu qu'on suppose $\pi > \pi_0$, et de plus $\pi^{\pi_0} = \pi$ (p. ex. $\pi = 2^c$ avec $c \geq \pi_0$).})

\begin{enumerate}
\item[a)] $I' \supset I'_0$.
\item[b)] $\card I' \leq \pi$.
\end{enumerate}

% ---- from Der-ChXIII-25-01-04_p046-050.tex
\printedpage{44}
% head: A. Grothendieck : Dérivateurs, Ch. XIII, Édition des Universités de Montpellier II et Paris VII

\begin{itemize}
\item[c)] \textit{$I'$ grand devant $\pi_0$.}
\end{itemize}

\texttt{[Le lemme tel qu'il est énoncé est évident pour les mêmes raisons que ci-dessus (même sans l'hypothèse de la note de bas de page). La démonstration qui suit prouve en fait un résultat plus fin :}

\texttt{\textsc{Lemme} 1$'$ : \textit{Soient $\pi_0$, $\pi$ deux cardinaux, avec $\pi_0<\pi$ et $\pi^{\pi_0}=\pi$, $I$ un ensemble ordonné grand devant $\pi_0$, et $I'_0$ une partie de $I$ de cardinal $\le\pi$. Alors il existe une partie $I'$ de $I$ satisfaisant les conditions :}}

\begin{itemize}
\item[\texttt{a)}] $I'\supset I'_0$\texttt{.}
\item[\texttt{b)}] $\card I'\le\pi$\texttt{.}
\item[\texttt{c)}] \texttt{\textit{$I'$ grand devant $\pi_0$. ]}}
\end{itemize}

On va essayer de construire $I'$ par récurrence transfinie, en définissant les $I'_\alpha\subset I$ ainsi :
\[
\begin{aligned}
I'_0 &= I'_0\;, \\
I'_\alpha &= \textstyle\bigcup_{\beta<\alpha} I'_\beta \quad \text{si $\alpha$ est un ordinal limite,} \\
I'_{\alpha+1} &= \left\{\parbox{0.6\textwidth}{déduit de $I'_\alpha$ ainsi : pour toute partie $J$ de $I'_\alpha$ telle que $\card J\le\pi_0$, qui n'est pas majorée dans $I'_\alpha$, on choisit un majorant $i_J$ dans $I$ (\texttt{[qui]} existe, car $I$ \texttt{[est]} grand devant $\pi$, donc devant $\pi_0$) et $I'_{\alpha+1}$ s'obtient en adjoignant ces $i_J$ à $I'_\alpha$.}\right.
\end{aligned}
\]

\msp{page 81}
Soit alors $\alpha_0$ le premier ordinal qui soit de cardinal $>\pi_0$, et considérons
\[
I' \;=\; I'_{\alpha_0} \;=\; \bigcup\limits_{\substack{\alpha<\alpha_0\\ \textit{i.e.}\ \card\alpha\le\pi_0}} I'_\alpha .
\]
Alors $I'$ contient $I'_0$, il est grand devant $\pi_0$, mais quel est son cardinal ? Prouvons par récurrence transfinie que
\[
\card I'_\alpha \;\le\pi \qquad\quad \text{si } \alpha<\alpha_0,\ \ \textit{i.e.}\ \ \card\alpha\le\pi_0 .
\]
C'est O.K. pour $\alpha=0$ par hypothèse sur $I'_0$. Si c'est vrai pour $I'_\alpha$, est-ce vrai pour $I'_{\alpha+1}$ ? On adjoint à $I'_\alpha$ un ensemble d'éléments $i_J$ qui est de cardinal $\le\card\mathfrak{P}_{\le\pi_0}(I'_\alpha)$, donc de cardinal $\le\card(I'_\alpha)^{\pi_0}\le\pi^{\pi_0}$, et si on suppose
\[
\pi^{\pi_0} \;=\; \pi,
\]
on gagne. Supposons enfin que $\alpha$ soit un ordinal limite (de cardinal $\le\pi_0$). Donc on a
\[
I'_\alpha \;=\; \bigcup_{\beta<\alpha} I'_\beta ,
\]
donc
\printedpage{45}
% head: A. Grothendieck : Dérivateurs, Ch. XIII, Édition des Universités de Montpellier II et Paris VII
\msp{page 82}
\[
\card I'_\alpha \;\le\; \underbrace{(\card\alpha)\pi}_{\le\pi\pi=\pi^2=\pi} \qquad\quad \text{puisque les } \card I'_\beta\le\pi,
\]
on gagne encore.

On trouve enfin, par le même argument,
\[
\card I'_{\alpha_0} \;\le\; \underbrace{(\card\alpha_0)\pi\le\pi^2}_{\substack{\text{si }\card\alpha_0\le\pi,\\ \textit{i.e.}\text{ si }\pi_0<\pi}} = \pi,
\]
donc on gagne encore.

\dgrblock{dgrb730}{Why the transfinite closure has to run far enough}
\begin{dcomment}{exposition}{p82:donc on gagne encore.}{Why the transfinite closure has to run far enough}
The editors distinguish two different closure problems. If $I$ is already grand devant $\pi$ and $S\subset I$ has size at most $\pi$, a single upper bound $s$ gives the required subposet $S\cup\{s\}$. It has a largest element, so it is grand devant every cardinal. This proves the stated Lemma~1 without any exponent assumption. The longer proof is useful when the ambient poset is only grand devant $\pi_0$; that is the corrected Lemma~$1'$.

Put $\lambda=\pi_0^+$. Starting with $S_0=S$, adjoin at each successor stage an upper bound for every subset of $S_\alpha$ having cardinal at most $\pi_0$, and take unions at limit stages. Suppose $\pi>\pi_0$ and $\pi^{\pi_0}=\pi$. Each successor step adds at most $\pi^{\pi_0}=\pi$ elements. A union of at most $\pi$ sets of size at most $\pi$ again has size at most $\pi$, so every $S_\alpha$, including $S_\lambda$, satisfies the desired bound.

The omitted verification is that $S_\lambda$ is grand devant $\pi_0$. Given $A\subset S_\lambda$ with $|A|\leq\pi_0$, choose a stage containing each of its elements. Since $\lambda=\pi_0^+$ is regular, these at most $\pi_0$ stages are bounded below $\lambda$. Hence $A\subset S_\alpha$ for some $\alpha<\lambda$, and the next stage contains an upper bound for $A$. This is why merely iterating the operation finitely many times, or through a sequence of cofinality at most $\pi_0$, would not prove the assertion.

The exponent hypothesis is sufficient rather than intrinsic. For example, when the desired bound is exactly $\pi_0^+$, enumerate the initial subset in a sequence of length $\pi_0^+$ and recursively choose an upper bound for the elements enumerated so far and the previously chosen bounds. At each proper stage there are at most $\pi_0$ requirements. Adjoining the resulting increasing sequence to the initial subset produces a subposet grand devant $\pi_0$, of cardinal at most $\pi_0^+$, without assuming $(\pi_0^+)^{\pi_0}=\pi_0^+$. We will use this elementary successor-cardinal improvement when returning to the unfinished final lemma.

These are two applications of the same method: control the number of requirements at a successor step, then choose a length whose cofinality catches every small family of requirements. Modern accessibility arguments still use this mechanism. The explicit stronger size assumption in \citet[Lemma~5.4.2.8]{Lurie2009HTT} serves exactly this purpose in the setting of $\infty$-categories.
\end{dcomment}


Ainsi, on a trouvé un deuxième système inductif
\[
(X_{\tilde\imath})_{\tilde\imath\in\tilde I}
\]
du même type que le système initial, ayant comme limite inductive $X$ (qu'il ait même limite inductive est immédiat). Mais qu'a-t-on gagné ? Ceci : dans $\tilde I$ toute partie de cardinal $\le\pi$ et grande devant $\pi_0$ est non seulement \textit{majorée} dans $\tilde I$, mais même y admet une borne supérieure. Cela résulte du % margin: double vertical rule beside the four lines "inductive est immédiat" ... "supérieure. Cela résulte du"

\msp{page 83}
\textbf{Lemme 2 :} \textit{Soit $I$ grand devant $\pi>\pi_0$, avec $\pi^{\pi_0}=\pi$} \texttt{[il suffit $\pi\ge\pi_0$]} \textit{(cf. lemme 1). Soit $(I'_\alpha)_{\alpha\in A}$ un ensemble de parties de $I$ qui sont dans $\tilde I$,} i.e.

\begin{itemize}
\item[a)] \textit{de cardinal $\le\pi$, et}
\item[b)] \textit{grandes devant $\pi_0$,}
\end{itemize}

\textit{et supposons que} $\card A\le\pi$ \textit{et que $A$, ordonné par inclusion, soit grand devant $\pi_0$. Alors $I'=\bigcup_{\alpha\in A} I'_\alpha$ est encore dans $\tilde I$,} i.e. \textit{$I'$ est de cardinal $\le\pi$ (car de cardinal $\le\pi^2=\pi$), et de plus, grande devant $\pi_0$.} (Vérification immédiate.)

Ainsi, on a prouvé ceci :

% margin: single vertical rule beside Proposition 1, from here to the end of the page
\textbf{Proposition 1.} \textit{Soit $E$} \texttt{[une]} \textit{grosse catégorie, avec filtration cardinale} $(\Filt^\pi)_{\pi\ge\pi_0}$. \textit{Soit $\pi$ un cardinal avec}
\[
\pi>\pi_0, \qquad\quad \pi^{\pi_0}=\pi \qquad\quad \textit{(p. ex. $\pi=2^c$, avec $c\ge\pi_0$)} \quad \texttt{[}\pi\ge\pi_0\ \texttt{suffit].}
\]
\textit{Alors pour tout $X$ dans $E$, on peut trouver un système inductif $(X_i)_{i\in I}$, de limite $X$, avec}

\begin{itemize}
\item[a)] \[ \left\{ \begin{aligned} & X_i\in\Filt^\pi \ \ \textit{pour tout $i\in I$,} \\ & \ I \ \textit{grand devant $\pi$,} \end{aligned} \right. \]
\end{itemize}

\textit{satisfaisant aux conditions supplémentaires suivantes :}

\begin{itemize}
\item[b)] \textit{Toute partie $J$ de $I$ de cardinal $\le\pi$, et grande devant $\pi_0$, admet une borne supérieure $i_J$ dans $I$.}
\item[c)] \textit{On a, sous ces conditions}
\end{itemize}

\printedpage{46}
% head: A. Grothendieck : Dérivateurs, Ch. XIII, Édition des Universités de Montpellier II et Paris VII
\msp{page 84}
\[
X_{i_J} \xleftarrow{\ \sim\ } \varinjlim_{i\in J} X_i .
\]
\textit{De plus, si $X\in\Filt^{\pi'}$ avec $\pi'>\pi$, on peut prendre $I$ de cardinal $\le\pi'^{\pi}$ (donc de cardinal $\le\pi'$ si $\pi'^{\pi}=\pi'$, p. ex. si $\pi'=2^{c'}$, avec $c'\ge\pi$).}

\textbf{Définition.} On appelle un système inductif $(X_i)_{i\in I}$, satisfaisant les conditions \textit{a}), \textit{b}), \textit{c}) ci-dessus, un \textit{système inductif bien adapté} (à la filtration cardinale donnée, pour les cardinaux $\pi_0$, $\pi$ donnés), et la représentation correspondante
\[
X \xleftarrow{\ \sim\ } \varinjlim X_i
\]
une \textit{représentation bien adaptée} de $X$ en termes d'objets de $\Filt^\pi$.

\dgrblock{dgrb731}{A presentation that is continuous on its small pieces}
\begin{dcomment}{exposition}{p84:représentation bien adaptée}{A presentation that is continuous on its small pieces}
Grothendieck wants the indexing set to remember small colimits already performed. An arbitrary directed presentation has upper bounds, but its value at an upper bound need not equal the colimit of the stages below it. Replacing indices by small directed subsets builds this continuity into the presentation.

Take $X=\varinjlim_{i\in I}X_i$ as in condition~(c), and let $P$ consist of the subposets $S\subset I$ of size at most $\pi$ which are grand devant $\pi_0$. Order $P$ by inclusion and put $Z_S=\varinjlim_{s\in S}X_s$. The ambient hypothesis on $I$ gives an upper bound for the union of at most $\pi$ members of $P$; adjoining that element produces another member of $P$. Thus $P$ is grand devant $\pi$. Condition~(b) gives $Z_S\in\Filt^\pi(E)$, while the elementary bound on subsets gives $|P|\leq|I|^\pi$.

The equality of the two colimits can be verified directly, without a cofinality convention. A compatible family of maps $Z_S\to Y$ gives maps $X_i\to Y$ by using the singleton subset $\{i\}$. For $i\leq j$, the two-element subset $\{i,j\}$, which has largest element $j$, ensures compatibility. Conversely a cocone on the $X_i$ induces one on every $Z_S$, and these two operations are inverse. Hence $\varinjlim_P Z_S\simeq X$.

Now let $A\subset P$ be grand devant $\pi_0$ and have cardinal at most $\pi$. Its union $S=\bigcup_{T\in A}T$ has size at most $\pi$. For a subset $B\subset S$ of size at most $\pi_0$, choose one $T_b\in A$ containing each $b$. An upper bound $T\in A$ for these members contains $B$, and an upper bound for $B$ inside $T$ belongs to $S$. Thus $S\in P$, and $S$ is the supremum of $A$. A cocone on the diagrams over all $T\in A$ is exactly a cocone over their union, so
\[
 Z_{\sup A}=Z_S\simeq\varinjlim_{T\in A}Z_T.
\]
This proves the continuity clause in Proposition~1, not merely the existence of suprema in its indexing poset.

When all filtered colimits exist and the filtration is strict in the bounded sense of footnote~41, the same proof uses ordinary directed subposets $S$ and ordinary directed families $A$. A union of a directed family of directed subposets is directed. A subset of $I$ of size at most $\pi$ can be enlarged to a directed subset of that same size by adding upper bounds for finite subsets through countably many stages. This yields Proposition~1 bis. The word ``perfectly'' records the stronger continuity, with no change to the cardinality bound on the partial diagrams. The inverse-limit sign in the preceding display on manuscript page~85 is a source slip: the construction uses a colimit.
\end{dcomment}


\textbf{NB} Quand $E$ est stable par petites $\varinjlim$ filtrantes, la situation se simplifie considérablement. On peut alors prendre
\[
\tilde I \;=\; \{I'\subset I \mid \card I'\le\pi,\ I' \text{ filtrant}\}
\]
(sans s'encombrer de la condition $I'$ grand devant $\pi_0$) et définir encore le système inductif $(X_{I'})_{I'\in\tilde I}$ par
\[
X_{I'} \;=\; \varinjlim_{i\in I'} X_i \;.
\]
On a encore
\msp{page 85}
\[
\card\tilde I \;\le\; (\card I)^\pi \;\le\; \pi'^{\pi} \qquad\quad \text{si }\ \card I\le\pi'^{\pi}
\]
(notamment, si $X\in\Filt^{\pi'}$, et qu'on prenne $I$ comme indiqué dans la condition $c$ des filtrations cardinales). De plus, cette fois, le fait que $\tilde I$ soit grand devant $\pi$ est trivial, on a même mieux, à part que $I$ soit filtrante :

\begin{itemize}
\item[$(*)$] \textit{Toute} partie $(I'_\alpha)_{\alpha\in A}$ de $\tilde I$ filtrante et de cardinal $\le\pi$ admet une borne supérieure $\tilde\imath_A$,
\end{itemize}

savoir la réunion (qui est bien filtrante). De plus, le système inductif $(X_{\tilde\imath})_{\tilde\imath\in\tilde I}$ a la propriété plus forte que
\[
X_{\tilde\imath_A} \xleftarrow{\ \sim\ } \varprojlim_{\tilde\imath\in A} X_{\tilde\imath} \;. % [sic] \varprojlim (arrow under lim pointing left) -> \varinjlim, inductive limit of the inductive system X_ĩ over the filtrant part A
\]
Ainsi on a la

\textbf{Proposition 1 bis.} \textit{Supposons $E$ stable par petites $\varinjlim$ filtrantes, et} $\Filt^\pi$ \textit{stable par celles-ci, où $\pi\ge\pi_0$ est donné. Alors pour tout $X$ dans $E$, on peut trouver un système inductif $(X_i)_{i\in I}$ de limite inductive $X$, avec}

\begin{itemize}
\item[a)] $X_i\in\Filt^\pi$ \textit{pour tout $i\in I$, $I$ ensemble ordonné grand devant $\pi$.}
\item[b)] \textit{Toute partie $J$ de $I$ filtrante de cardinal $\le\pi$ admet une borne supérieure $i_J$ dans $I$.}
\item[c)] \textit{Sous ces conditions}
\end{itemize}

\printedpage{47}
% head: A. Grothendieck : Dérivateurs, Ch. XIII, Édition des Universités de Montpellier II et Paris VII
\msp{page 86}
\[
X_{i_J} \;=\; \varinjlim_{i\in J} X_i \;.
\]
De tels systèmes inductifs seront dits \textit{“parfaitement adaptés”} (à la filtration cardinale donnée, pour le cardinal $\pi$ donné – $\pi_0$ n'y figure plus, si ce n'est comme début de la filtration cardinale sans plus).

\dgrblock{dgrb732}{Sets and modules as perfectly adapted presentations}
\begin{dcomment}{examples}{p86:De tels systèmes inductifs seront dits}{Sets and modules as perfectly adapted presentations}
The basic example makes every condition concrete. In $E=\mathrm{Set}$ put $\Filt^\pi(E)=\{S:|S|\leq\pi\}$. A skeleton is a set of cardinals, so the stage is essentially small. A colimit of a diagram with at most $\pi$ objects, each of size at most $\pi$, is a quotient of a disjoint union of size at most $\pi\cdot\pi=\pi$. Thus the required bounded closure holds, even for arbitrary colimits of that size.

For a set $X$ consider
\[
 I=[X]^{\leq\pi},\qquad X_A=A\quad(A\in I),
\]
with inclusion maps. The union of at most $\pi$ members has size at most $\pi$, so $I$ is grand devant $\pi$. Every nonempty directed family $J\subset I$ of size at most $\pi$ has supremum $\bigcup J$, and its colimit in sets is this same union. The colimit of the entire diagram is $X$, since the singletons cover it. If $|X|\leq\pi'$, then $|I|\leq(\pi')^\pi$ for infinite $X$; finite $X$ causes no difficulty. This is a perfectly adapted representation with the precise bound appearing in the source.

A set $A$ of size at most $\pi$ is $\pi^+$-presentable. To factor a map $A\to\varinjlim_IY_i$ through one stage, choose representatives for its at most $\pi$ values, and then one upper bound for their indices. If two such factorizations induce the same map, each of their at most $\pi$ component equalities holds at some later stage; another upper bound makes the maps equal. Both surjectivity and injectivity of the Hom comparison are used. They are the categorical counterparts of choosing generators and then checking relations.

For modules over a ring $R$, fix an infinite $\pi\geq|R|$. Use modules whose underlying sets have size at most $\pi$. They form an essentially small class, and bounded filtered colimits preserve the size bound. Every module $X$ is the union of its submodules of size at most $\pi$. Any set of at most $\pi$ elements generates such a submodule: a linear combination uses finitely many coefficients, so there are at most $\pi$ combinations. The poset of these submodules is grand devant $\pi$, and directed unions of at most $\pi$ of them again have size at most $\pi$. This gives the same perfectly adapted presentation.

The module presentability assertion also includes relations. A module of size at most $\pi$ has a surjection from a free module on at most $\pi$ generators. That free module and its kernel have size at most $\pi$, so the kernel is generated by at most $\pi$ elements. Factoring a homomorphism through a $\pi^+$-directed colimit requires at most $\pi$ choices for generators and at most $\pi$ equations for relations; a common later stage handles them all. Equality of two factorizations is checked on the generators. Thus the chosen stage consists of $\pi^+$-presentable modules.

These examples also locate an omission in the glossary. There, ``strict filtration'' is stated as closure of each stage under all filtered colimits, with no bound on the indexing cardinal. Together with condition~(c), that statement would force $E=\Filt^\pi(E)$ for every $\pi$, hence would make $E$ essentially small. The definition actually used here is the bounded one in footnote~41: closure under filtered diagrams of cardinal at most $\pi$. For instance, the union of $\pi^+$ many successively larger subsets of a set of cardinal $\pi^+$ shows why the bound cannot be discarded.
\end{dcomment}


Comme pratiquement, on aura à manier surtout des $E$ stables par petites $\varinjlim$ filtrantes, et des filtrations cardinales stables en conséquence (renforcement de la condition \textit{b}) des filtrations cardinales), c'est le cas, et la notion de \textit{représentation parfaitement adaptée} de $X$ comme $\varinjlim\limits_{I} X_i$ ($X_i\in\Filt^\pi$), que je mettrai en avant. Le cas “général”, où on suppose seulement \texttt{[que]} $E$ satisfait $L_{\pi_0}$, est plus alambiqué et se traite moins joliment (\footnote[41]{On parlera de filtration cardinale \textit{stricte,} si les $\underline{\Filt}^\pi$ sont stables par $\varinjlim$ filtrantes de cardinal $\le\pi$ (et que ces limites filtrantes petites existent dans $E$).}).

\textbf{Définition 2.} Soit $E$ munie d'une filtration cardinale stricte $(\underline{\Filt}^\pi)_{\pi\ge\pi_0}$. Soit $M$ une partie de $\Ob E$, $\underline{M}$ la sous-catégorie pleine de $E$ qu'elle définit, et soit $\pi$ un cardinal $\ge\pi_0$. On dit que $M$ est \textit{$\pi$-accessible,} où que $\underline{M}$ est \textit{$\pi$-accessible}, % [sic] où que -> ou que, disjunction "ou" not the relative "où"
\msp{page 87}
si les conditions suivantes sont satisfaites. Pour tout système inductif
\[
(X_i)_{i\in I}
\]
parfaitement adapté à la filtration cardinale pour $\pi$ (donc $I$ ensemble ordonné filtrant tel que toute partie filtrante $J$ de cardinal $\le\pi$ ait une borne supérieure $i_J$, et $X_{i_J}\xleftarrow{\ \sim\ }\varinjlim\limits_{j\in J} X_j$) ; soit $I_M$ l'ensemble des $i\in I$ tels que $X_i\in M$. Ceci posé, on veut ceci : pour que $X$ ($=\varinjlim X_i$) soit dans $M$, il faut et il suffit que $I_M$ soit cofinal dans $I$ (de sorte que $X$ apparaît comme $\varinjlim$ filtrante, avec ensemble d'indices $I_M$ grand devant $\pi$, d'éléments de $M\cap\underline{\Filt}^\pi\overset{\text{déf}}{=}M^\pi$).

\msp{page 88}
J'ai envie de reprendre la définition de la $\pi$-accessibilité de $M$, pour une filtration cardinale donnée, à coups de conditions équivalentes. Pour tout système inductif $(X_i)_{i\in I}$ dans $M$ \texttt{[plutôt dans $E$]}, je garde la notation
\[
I_M \;=\; \{i\in I \mid X_i\in M\}.
\]
Je pose, pour tout cardinal $\pi\ge\pi_0$,
\[
M^\pi \;=\; M\cap\Filt^\pi .
\]

\begin{itemize}
\item[\circled{a$_\pi$}] $M$ est l'ensemble des objets de la forme
\[
\varinjlim_{i\in I} X_i \;,
\]
où les $X_i$ sont dans $M^\pi$, et $I$ est grand devant $\pi$.
\end{itemize}

\printedpage{48}
% head: A. Grothendieck : Dérivateurs, Ch. XIII, Édition des Universités de Montpellier II et Paris VII

\begin{itemize}
\item[\circled{b$_\pi$}] Soit $(X_i)_{i\in I}$ un système inductif dans $M$ \texttt{[plutôt dans $E$]} parfaitement adapté à la filtration cardinale pour $\pi$, \textit{i.e.} $I$ filtrante et toute partie filtrante $J$ de $I$ de cardinal $\le\pi$ admet une borne supérieure $i_J$, et $X_{i_J}\xleftarrow{\ \sim\ }\varinjlim\limits_{J} X_i$, et enfin $X_i\in\Filt^\pi$ pour tout $i\in I$. Soit $X=\varinjlim X_i$. Pour que l'on ait $X\in M$, il faut et il suffit que $I_M$ soit cofinal dans $I$.
\item[\circled{c$_\pi$}] $\underline{M}^\pi=\underline{M}\cap\underline{\Filt}^\pi$ est une sous-catégorie de $\underline{M}$ génératrice par épimorphismes stricts, $\underline{M}$ est stable par $\varinjlim$ grandes devant $\pi$, enfin pour tout $X\in M$, $\underline{M}^\pi/X$ est une catégorie filtrante.
\end{itemize}

\msp{page 89}
Je présume que ces conditions sont équivalentes.

\dgrblock{dgrb733}{Three different ways to recognize a part}
\begin{dcomment}{exposition}{p89:Je présume que ces conditions sont équivalentes.}{Three different ways to recognize a part}
The purpose of the new definition is to recognize membership in a class $M$ from small approximations. This would be particularly useful when $E$ is an arrow category and $M$ is a class of weak equivalences. Grothendieck proposes three tests, but their logical strengths differ.

The first, $a_\pi$, says that $M$ is exactly the class of $\pi^+$-directed colimits of objects of $M^\pi=M\cap\Filt^\pi(E)$. It includes a closure assertion as well as an existence assertion. The second, $b_\pi$, says that \emph{every} perfectly adapted presentation detects membership by a cofinal set of $M$-stages. It therefore imposes a condition on the individual stages of a chosen presentation, beyond the existence of some good presentation. The third, $c_\pi$, combines strict generation, closure under $\pi^+$-directed colimits, and ordinary filteredness of $M^\pi/X$.

At a useful cardinal, with $M^\pi$ consisting of $\pi^+$-presentable objects, $a_\pi$ is the familiar description by an $\operatorname{Ind}_{\pi^+}$-completion. Indeed the canonical colimit functor out of that completion is fully faithful; the Hom calculation is the one in \citet[I, 9.18(a)]{SGA4I1972}. It follows in particular that the resulting full subcategory is closed under $\pi^+$-directed colimits. This explains the modern accessibility of the inclusion, but it does not make $b_\pi$ automatic.

There are two distinct issues in the proposed equivalence. Ordinary filteredness controls finite diagrams, whereas $a_\pi$ asks for compatibility of up to $\pi$ constraints at once. Separately, two cofinal collections of small approximations need not meet: one may pick only successor stages and another only limit stages. Continuity of the second presentation does not force membership to pass to its countable colimits. The following examples isolate these two failures while keeping all objects in the small stage genuinely presentable.
\end{dcomment}


\dgrblock{dgrb734}{Two counterexamples to the proposed equivalence}
\begin{dcomment}{development}{p89:Je présume que ces conditions sont équivalentes.}{Two counterexamples to the proposed equivalence}
The distinction just made can be checked in a category with no nontrivial parallel arrows. This removes any possible ambiguity about coherence and leaves the cardinal issue exposed. We give the examples as our own verification of the conjecture on page~89. We did not find these particular examples in the relevant portions inspected of SGA~4 I, \S9, \emph{Les D\'erivateurs} XVIII, \S4.13 and Lemma~4.16.7.11, and the accessible-category references listed below; no priority is asserted.

\textbf{Proposition: the implication $a_\pi\Rightarrow b_\pi$ fails at the same cardinal.}\ours
Take $\pi=\aleph_0$ and let $E$ be the poset $\mathcal P(\omega_1)$, viewed as a category under inclusion. Give it the strict filtration by cardinality, $\Filt^\rho(E)=\{S\subset\omega_1:|S|\leq\rho\}$. Let
\[
 M=\{\alpha+1:\alpha<\omega_1\}\cup\{\omega_1\},
\]
where ordinals denote their sets of predecessors. Then $a_\pi$ and $c_\pi$ hold, while $b_\pi$ fails. Every object of $\Filt^\pi(E)$ is $\pi^+$-presentable in $E$.

\emph{Proof.} Colimits in $E$ are unions. Each filtration stage is small, since it is a subset of the universe-small set $\mathcal P(\omega_1)$. A union of at most $\rho$ sets of size at most $\rho$ again has size at most $\rho$, proving the bounded closure axiom for every infinite $\rho$. Each $S\subset\omega_1$ is the union of its subsets of size at most $\rho$; their poset is $\rho^+$-directed because a union of at most $\rho$ such subsets is another one. If $|S|\leq\rho'$ and $\rho'>\rho$, the number of indices is at most $(\rho')^\rho$; the finite cases satisfy this bound as well. This verifies the full filtration axiom~(c), including its size estimate. Finally, in a $\pi^+$-directed union every countable subset of the union is contained in one stage. Since Hom sets in $E$ have at most one element, this factorization assertion proves presentability of all the countable objects.

The objects of $M^\pi$ are exactly the countable successor ordinals. Each is a constant directed colimit of such objects, and $\omega_1$ is the union of all countable successor ordinals; their ordered set is countably directed. Conversely, consider a countably directed family of countable successor ordinals. If its union $\beta$ is countable and is not attained, choose a countable cofinal subfamily. A single member bounding this subfamily would contain its union, contradicting nonattainment. Therefore a countable union arising this way is an attained successor ordinal. An uncountable union is $\omega_1$. This proves $a_\pi$ and closure of $M$ under countably directed colimits.

For $X\in M$, the slice $M^\pi/X$ is a nonempty linearly ordered set, hence is filtered. Its union is $X$. In a poset this says exactly that the canonical family of arrows from $M^\pi$ is strictly epimorphic: any common upper bound of that family in $M$ is an upper bound of $X$, and the required factorization is unique. Thus $c_\pi$ holds as well.

Now index a diagram by the nonzero countable limit ordinals $\delta<\omega_1$, and set $X_\delta=\delta$. The index is countably directed: the supremum of any countable family is still countable, and a larger countable limit ordinal is an upper bound. Every nonempty directed subset of this index of cardinal at most $\aleph_0$ has its supremum in the index; a supremum of nonzero limit ordinals is again a nonzero limit ordinal. At that supremum the diagram value is the union of the earlier values. The presentation is therefore perfectly adapted, and its colimit is $\omega_1\in M$. Nevertheless no stage belongs to $M$. Its set $I_M$ is empty. This disproves $b_\pi$.\hfill$\square$

\textbf{Proposition: the implication $c_\pi\Rightarrow a_\pi$ also fails.}\ours
In the same filtered category $E$, replace $M$ by
\[
 N=\{S\subset\omega_1:S\text{ is finite}\}\cup\{\omega_1\}.
\]
Then $c_\pi$ holds for $N$, but $a_\pi$ does not.

\emph{Proof.} Here $N^\pi$ consists of the finite subsets. For a finite $X$, the slice $N^\pi/X$ has largest element $X$; for $X=\omega_1$, it is the directed set of finite subsets, whose union is $\omega_1$. These unions prove strict generation, as in the first example. A countably directed family of finite subsets has finite union. Otherwise choose countably many distinct elements of the union and a stage containing each; an upper bound for these countably many stages would be a finite set containing infinitely many elements. A countably directed diagram in $N$ that contains $\omega_1$ has union $\omega_1$. Thus $N$ is closed under the required colimits, proving $c_\pi$.

But no countably directed colimit of objects of $N^\pi$ can equal $\omega_1$, by the same finite-union argument. Hence $a_\pi$ fails.\hfill$\square$

The examples are legitimate $\mathfrak U$-categories. If one requires the ambient category itself to be non-small, take $E\times\mathrm{Set}$ with
\[
 \Filt^\rho(E\times\mathrm{Set})=
 \{(S,T): |S|\leq\rho,\ |T|\leq\rho\}.
\]
A small skeleton is obtained by retaining the actual subsets $S\subset\omega_1$ and replacing $T$ by a cardinal. Colimits are computed in each coordinate, so bounded closure follows from the union calculation above and the set-colimit calculation in the block ``Sets and modules as perfectly adapted presentations.'' For condition~(c), use the product poset $[S]^{\leq\rho}\times[T]^{\leq\rho}$. It is $\rho^+$-directed and has colimit $(S,T)$. If both coordinates have size at most $\rho'>\rho$, its cardinal is at most $(\rho')^\rho\cdot(\rho')^\rho=(\rho')^\rho$. Countable objects are $\aleph_1$-presentable by the same simultaneous choice of representatives and equations in the two coordinates. Restrict the second coordinate of $M$ or $N$ to the empty set. All presentations, strict generating families and counterexample diagrams used above then have empty second coordinate, so every preceding verification remains valid. The repair to $c_\pi$ is substantial: one needs a $\pi^+$-filtered canonical slice with the asserted canonical colimit, rather than merely an ordinary filtered slice. The first example shows that even this repaired generation condition does not recover the universal-presentation condition $b_\pi$ at the same rank.
\end{dcomment}


$b_\pi\Longrightarrow a_\pi$. On suppose \textit{b}) vérifié, on veut prouver \textit{a}). D'abord que tout $X\in M$ est de la forme dite $\varinjlim\limits_{I} X_i$, avec les $X_i\in M^\pi$ et $I$ grand devant $\pi$. Par \texttt{[la]} proposition 1, on sait qu'on a une présentation parfaite adaptée pour $\pi$, $X=\varinjlim\limits_{i\in I} X_i$. Par \textit{b}) “il faut”, la partie $I_M$ de $I$ est cofinale dans $I$, donc comme $I$ est grand devant $\pi$, $I_M$ aussi, et on aura (posant $J=I_M$)
\[
X \;=\; \varinjlim_{i\in J} X_i, \qquad\quad X_i\in M^\pi,\ \ J \text{ grand devant } \pi,
\]
O.K. Inversement, prouvons que si $X$ se met sous la forme dite dans \textit{a}), il est dans $M$. Cela va résulter du

\textbf{Lemme 1 :} \textit{Si la condition} b) (\textit{ou} a)) \textit{est satisfaite, $M$ est stable par limites inductives grandes devant $\pi$.}

Il faut prouver que si
\msp{page 90}
\[
X \;=\; \varinjlim_{i\in I} X_i \;,
\]
avec $I$ grand devant $\pi$, les $X_i$ dans $M$, alors $X\in M$. La condition \textit{b}) “il suffit” l'assure quand on suppose de plus les $X_i$ dans $\Filt^\pi$ (donc dans $M^\pi$), de sorte que l'on aura $I_M=I$, donc $I_M$ cofinal, donc $X\in M$ (\footnote[42]{Pas tout à fait – on y suppose plus, sur le système inductif. On ne peut se dispenser, même dans le cas $X_i\in M^\pi$ $\forall\, i\in I$, faire le raisonnement des pages suivantes.}). Si on ne suppose pas les $X_i$ dans $\Filt^\pi$, il faut travailler un peu plus.

Passer tout de suite à la partie de la démonstration ci-dessous, page 96, qui prouve \textit{b}) $\Longrightarrow$ \textit{a}). Le reste de la démonstration va prouver le

\textbf{Lemme 2.} \textit{Si tout élément de Filt$^\pi$ est $\pi$-accessible, alors} b$_\pi$) \textit{(et même} a$_\pi$\textit{) implique que $M$ est stable par $\varinjlim$ grandes devant $\pi$.}

% ---- from Der-ChXIII-25-01-04_p051-055.tex
\printedpage{49}
% head: A. Grothendieck : Dérivateurs, Ch. XIII, Édition des Universités de Montpellier II et Paris VII
% foot: Version 25 janvier 2004

\msp{page 91}Soit
\[ \mathcal{C} \;=\; \underline{M}^\pi/\mathcal{X} \]
la catégorie suivante. Ses objets sont les triples
\[ (Z,i,\alpha) \quad \text{tels que} \quad Z\in M^\pi,\ i\in I,\ \alpha : Z\longrightarrow X_i \ , \]
les flèches
\[ (Z,i,\alpha) \;\longrightarrow\; (Z',i',\alpha') \]
sont les couples $(u,\lambda)$ avec
\[ u : Z\longrightarrow Z', \qquad \lambda : i\longrightarrow i', \qquad (\textit{i.e.}\ \text{on a } i\le i' \text{ et } \lambda \text{ est l'unique flèche de } i \text{ dans } i'), \]
tels que le diagramme
\[
\begin{tikzcd}[row sep=large, column sep=large]
Z \arrow[r,"\alpha"] \arrow[d,"u"'] & X_i \arrow[d,"\text{transition } \lambda_{i',i}"] \\
Z' \arrow[r,"\alpha'"] & X_{i'}
\end{tikzcd}
\]
commute. On peut donc définir $\mathcal{C}$ comme le produit fibré dans le carré cartésien
\[
\begin{tikzcd}[row sep=large, column sep=large]
 & \mathcal{C} \arrow[r,"\text{cofibré}"] \arrow[d] \arrow[dl,dotted,"\mathcal{Z}"'] & I \arrow[d,"\mathcal{X} : X \longmapsto X_i"] \\ % [sic] X \longmapsto X_i -> i \longmapsto X_i, the functor \mathcal{X} goes from I, so the variable should be i
\underline{M}^\pi & \underline{M}^\pi/E \arrow[l,dotted,"\text{source}"'] \arrow[r,"\text{but}","\text{cofibré}"'] & E.
\end{tikzcd}
\]
La fibre $\mathcal{C}_i$ de $\mathcal{C}$ en $i\in I$ est $\mathrm{Filt}^\pi/X_i$ \texttt{[plutôt} $M^\pi/X_i$\texttt{]}. On a, comme \textit{b}) $\Longrightarrow$ première moitié de \textit{a}), que $\underline{M}^\pi$ est un
\msp{page 92}système de générateurs stricts de $\underline{M}$ \texttt{[cela n'est pas clair sans autres hypothèses sur} $\underline{M}$\texttt{]}, et $X_i\in\operatorname{Ob}\underline{M}$ par hypothèse,
\[ X_i \;=\; \varinjlim_{Z\in\operatorname{Ob}\underline{M}^\pi/X_i} Z \;=\; \varinjlim_{\mathcal{C}_i}\mathcal{Z}|\mathcal{C}_i, \]
donc le système inductif donné $(X_i)_{i\in I}$ sur $I$ à valeurs dans $X$, se déduit du système inductif (non ordonné en général) sur $\mathcal{C}$, à valeurs dans $\underline{M}^\pi$, en ``intégrant sur les fibres'' de $\mathcal{C}\longrightarrow I$. D'ailleurs, comme $\underline{M}^\pi/E\longrightarrow E$ est cofibrée, donc propre, il en est de même de $\mathcal{C}\longrightarrow I$, donc les $\varinjlim$ de $\mathcal{Z}$ sur les fibres $\mathcal{C}_i$ ($=$ $\underline{M}^\pi/X_i$) forment bien un foncteur covariant $I\longrightarrow E$, dont la limite inductive est celle de $\mathcal{Z}$ (\footnote[43]{\textbf{NB} Si on admet la validité de $a$), alors il en résulte déjà que $\varinjlim\limits_{i\in I} X_i \simeq \varinjlim\limits_{\alpha\in C} Z_\alpha$ est dans $M$. (Stabilité de $M$ par $\varinjlim$ grandes devant $\pi$).}).

\printedpage{50}
% head: A. Grothendieck : Dérivateurs, Ch. XIII, Édition des Universités de Montpellier II et Paris VII
% foot: Version 25 janvier 2004

\msp{page 93}Soit, pour tout $i\in I$, $\tilde{\mathcal{C}}_i$ l'ensemble ordonné
\[ \tilde{\mathcal{C}}_i = \{\text{sous-catégories } \mathcal{C}'_i \text{ de } \mathcal{C}_i = \underline{M}^\pi/X_i \mid \operatorname{card}\mathcal{Z}(\mathcal{C}'_i)\le\pi,\ \ \mathcal{C}'_i \text{ filtrante}\}. \]
J'admets pour l'instant le

\textbf{Lemme} (\footnote[44]{Utilise l'hypothèse tacite que les éléments de $\mathrm{Filt}^\pi$ sont $\pi$-accessibles.}). \textit{Pour tout $X\in M$, la catégorie $\underline{M}^\pi/X$ est filtrante. (Sous l'hypothèse} b), \textit{ou que tout objet de $M$ est de la forme dite dans} a).\textit{)}

\dgrblock{dgrb735}{Presentability supplies the strict generators}
\begin{dcomment}{exposition}{p93:est filtrante. (Sous l'hypothèse}{Presentability supplies the strict generators}
Grothendieck now tries to replace a diagram of large $M$-objects by one diagram of small $M$-objects. The comma category is the right construction: it records a small object together with the map by which it participates in a particular large object. The editors' reservation on page~92 concerns the claim that a chosen colimit presentation already makes \emph{all} maps from small objects into a canonical colimit presentation. Presentability supplies the missing comparison.

Here is the precise assertion. Let $\kappa=\pi^+$, let $A\subset E$ be an essentially small full subcategory of $\kappa$-presentable objects, and suppose
\[
 X\simeq\varinjlim_{j\in J}Y_j,
 \qquad J\ \kappa\text{-directed},\qquad Y_j\in A.
\]
Then $A/X$ is $\kappa$-filtered, and its canonical diagram has colimit $X$. This is the canonical-diagram form of the fully faithful $\operatorname{Ind}_\kappa$ construction in \citet[I, 9.18(a)]{SGA4I1972}.

To prove the assertion, take any diagram of fewer than $\kappa$ morphisms in $A/X$. Each of its objects is a map $T\to X$ with $T$ $\kappa$-presentable, so it factors through a $Y_j$. Choose one common upper bound for the indices used. For every arrow in the given diagram, the two resulting maps into this common stage agree after passage to $X$. Presentability of its source makes that equality hold at some later stage. There are fewer than $\kappa$ equalities, so another common upper bound makes them all hold together. The resulting object $Y_{j'}\to X$ is a cocone in $A/X$. This proves $\kappa$-filteredness, including the case of the empty diagram.

To identify the colimit, let a compatible family $T\to Z$ be given for all objects $(T\to X)$ of $A/X$. Restrict it to the chosen presentation $Y_j\to X$, producing a unique map $h:X\to Z$. Any $T\to X$ factors through a $Y_j$, and compatibility with that factorization says that the prescribed map $T\to Z$ is $h$ composed with $T\to X$. Thus $h$ induces the whole family, uniquely. This proves the universal property of the canonical colimit. In particular its canonical family is strictly epimorphic in the ambient category, hence also in any full subcategory containing $X$ and $A$.

Apply this with $A=M^\pi$. The first half of $a_\pi$, already obtained from $b_\pi$, supplies the required presentations. The tacit hypothesis in footnote~44 supplies presentability. Under those hypotheses the lemma on page~93 is correct, and in fact its slice is $\pi^+$-filtered, a stronger conclusion than printed. The precise missing step is factorization and eventual equality for \emph{arbitrary} maps from $M^\pi$, not an additional assumption about epimorphisms in $M$.

The same argument explains why this construction is stable under a sufficiently directed change of the large object. For a $\kappa$-directed diagram $(X_i)$ in $M$, form the category $C$ of triples $(T,i,T\to X_i)$. A diagram of fewer than $\kappa$ arrows in $C$ can first be moved into one fibre over $I$, and then given a cocone in that fibre by the preceding paragraph. Thus $C$ is $\kappa$-filtered. Its colimit is $\varinjlim_I X_i$, because giving a cocone on $C$ is giving compatible cocones on all the fibres. After the usual cofinal directed-poset replacement, this is a presentation by objects of $M^\pi$. If $a_\pi$ holds, it proves closure of $M$ under $\kappa$-directed colimits. The directed-poset replacement is the ordinary-category case of \citet[Proposition~5.3.1.18]{Lurie2009HTT}; it changes the index, not the colimit.
\end{dcomment}


Il en résulte, comme $X_i\in M$, que $\mathcal{C}_i$ est filtrante, donc la réunion des $\mathcal{C}'_i\in\tilde{\mathcal{C}}_i$ est $\mathcal{C}_i$ ; l'ensemble ordonné $\tilde{\mathcal{C}}_i$ est filtrant, et toute partie $(\mathcal{C}'_{i_\alpha})_{\alpha\in A}$ filtrante et de cardinal $\le\pi$ de $\tilde{\mathcal{C}}_i$ y admet une borne supérieure $\mathcal{C}'_{i_A}$, savoir la réunion des $\mathcal{C}'_{i_\alpha}$. $\tilde{\mathcal{C}}_i$ est donc un ensemble ordonné $\pi$-adapté.

On peut considérer les ensembles ordonnés $\tilde{\mathcal{C}}_i$, pour $i\in I$ variable, comme les fibres d'une application croissante d'ensembles ordonnés
\[ p \;:\; \tilde{\mathcal{C}} \;\longrightarrow\; \mathcal{C}, \]
en écrivant
\[ \left\{ \begin{array}{l}
(\mathcal{C}'_i\subset\mathcal{C}_i) \;\le\; (\mathcal{C}'_j\subset\mathcal{C}_j) \\
\text{si et seulement si } i\le j, \\
\text{et le foncteur canonique de transition } \mathcal{C}_i\longrightarrow\mathcal{C}_j \text{ applique } \mathcal{C}'_i \text{ dans } \mathcal{C}'_j.
\end{array}\right. \]
On voit donc que $\tilde{\mathcal{C}}$ est lui aussi filtrant et grand devant $\pi$.

\msp{page 94}On définit un système inductif sur $\tilde{\mathcal{C}}$, à valeurs dans $\underline{\mathrm{Filt}}^\pi$, en posant
\[ Z_{\mathcal{C}'_i} \;=\; \varinjlim_{\alpha\,\in\,\mathcal{C}'_i\ (\subset\,\mathcal{C}_i\,=\,\underline{M}^\pi/X_i)} Z_\alpha \]
(où un $\alpha\in\mathcal{C}_i$ est formé, je le rappelle, de $Z_\alpha\xrightarrow{u_\alpha}X_i$, avec $Z_\alpha\in M^\pi\subset\underline{\mathrm{Filt}}^\pi$). Ces $Z_{\mathcal{C}'_i}$ ne sont plus forcément dans $M^\pi$, mais seulement dans $\mathrm{Filt}^\pi$. (Alors à quoi ça a servi de définir les $\mathcal{C}_i$ eux-mêmes comme $\underline{M}^\pi/X_i$, au lieu de $\underline{\mathrm{Filt}}^\pi/X_i$ ?)

Le système inductif
\[ (Z_{\mathcal{C}'_i})_{\mathcal{C}'_i\in\tilde{\mathcal{C}}_i} \]
est parfaitement bien adapté à la filtration cardinale pour le cardinal $\pi$, et il est immédiat que sa $\varinjlim$ est isomorphe à $\varinjlim\limits_{\alpha\in\operatorname{Ob}\mathcal{C}_i=\underline{M}^\pi/X_i} Z_\alpha = X_i\in M$ . Donc par \textit{b}), l'ensemble $\tilde{\mathcal{C}}_M$ des $\mathcal{C}'_i\in\tilde{\mathcal{C}}_i$ tels que $Z_{\mathcal{C}'_i}\in M$ (donc $Z_{\mathcal{C}'_i}\in M^\pi$) est cofinal dans $\tilde{\mathcal{C}}_i$. Il s'ensuit donc que l'ensemble $\tilde{\mathcal{C}}_M$ est lui aussi cofinal dans $\tilde{\mathcal{C}}$.

\printedpage{51}
% head: A. Grothendieck : Dérivateurs, Ch. XIII, Édition des Universités de Montpellier II et Paris VII
% foot: Version 25 janvier 2004

\msp{page 95}On vérifie aussi immédiatement que
\[ \varinjlim_{\alpha\in\tilde{\mathcal{C}}} Z_\alpha \;=\; \varinjlim X_i \;=\; X, \]
donc, $\tilde{\mathcal{C}}_M$ étant cofinal,
\[ \varinjlim_{\alpha\in\tilde{\mathcal{C}}_M} Z_\alpha \;=\; X. \]
Au total, on est arrivé à représenter $X$ sous la forme
\[ X \;=\; \varinjlim_{\alpha\in A} X_\alpha, \]
où cette fois les $X_\alpha$ sont \textit{dans} $M^\pi$ (pas seulement dans $M$), et $A$ est grand devant $\pi$. Mais on ne sait pas que toute partie filtrante $B$ de cardinal $\le\pi$ de $A$ admet une borne supérieure $\alpha_B$ dans $A$, et encore moins que $X_{\alpha_B} = \varinjlim\limits_{\alpha\in B} X_\alpha$ (\footnote[45]{Si on suppose $(a_\pi)$, cet argument prouve du moins que $M$ est stable par $\varinjlim$ grandes devant $\pi$.}). Donc on ne peut encore appliquer directement \textit{b}). S'il ne s'agissait que de prouver la deuxième moitié de \textit{a}), où on suppose déjà les $X_i\in M^\pi$ et $I$ grand devant $\pi$, on n'a rien gagné du tout par nos labeurs !

\dgrblock{dgrb736}{The size of an indexing category cannot be replaced by the size of its image}
\begin{dcomment}{development}{p95:Donc on ne peut encore appliquer directement}{The size of an indexing category cannot be replaced by the size of its image}
The partial-colimit construction is meant to keep all intermediate objects within $\Filt^\pi$. For that purpose it must bound the indexing category. The literal condition on page~93, $\operatorname{card}\mathcal Z(C_i')\leq\pi$, bounds only the image of the diagram. That is insufficient, even if its cardinal counts image morphisms as well as image objects.

\textbf{A filtered counterexample to the image bound.}\ours
Let $E=M=\mathrm{Set}$ with the usual cardinal filtration, let $\pi=\aleph_0$, and choose an uncountable set $X$. Inside $M^\pi/X$, take the category $C'$ whose objects are injections $a:n\hookrightarrow X$ from the standard finite sets $n=\{0,\ldots,n-1\}$. A morphism from $a:n\hookrightarrow X$ to $b:m\hookrightarrow X$ is a map $u:n\to m$ satisfying $bu=a$.

This category is filtered. Two objects map into an injection whose image is the union of their two finite images. Parallel arrows are equal, since their composites with the injective target map are equal. The source functor $\mathcal Z:C'\to\mathrm{Set}$ has image inside the countable category of standard finite sets and maps; hence $\operatorname{card}\mathcal Z(C')\leq\aleph_0$. Yet
\[
 \varinjlim_{C'}\mathcal Z\simeq X.
\]
Indeed every element of $X$ occurs through a one-element injection, and two representatives determine the same element precisely when they become equal in a finite union of their images. Equivalently the compatible injections give the usual directed union of the finite subsets of $X$. This colimit is not in $\Filt^{\aleph_0}(E)$.\hfill$\square$

The corrected construction uses filtered \emph{indexing categories} $D\subset C_i$ with at most $\pi$ morphisms, or directed indexing posets of that size mapping to $C_i$. This correction does not require a bound on every Hom set of $C_i$. The subcategories need not be full. Start with at most $\pi$ specified objects and arrows; close under identities and composition, adjoin cones for finite diagrams, and repeat countably often. Each step adds at most $\pi$ data, since $\pi$ is infinite. The resulting filtered subcategory still has at most $\pi$ morphisms. In particular the poset of these small filtered subcategories is grand devant $\pi$: perform this construction on the union of any at most $\pi$ of them.

A directed union of at most $\pi$ such subcategories is again a filtered subcategory of size at most $\pi$ and is their supremum. Colimits over this union are the colimits of the partial colimits. To apply the source's filtration axiom, which is phrased for posets, a filtered category with at most $\pi$ arrows can be replaced cofinally by a directed poset of size at most $\pi$: the usual finite-diagram construction adjoins finitely many compatibility data at each of countably many stages. Thus the corrected partial colimit lies in $\Filt^\pi$, and its presentation over the poset of small subcategories is perfectly adapted.

This also repairs the rest of the argument on pages~93--95. The map from the total poset of pairs $(i,D\subset C_i)$ has target $I$, as its displayed order relation specifies; the printed target $C$ is a slip. To bound at most $\pi$ such pairs, move their indices to one index of $I$ and enlarge the union of their images to a small filtered subcategory of the corresponding fibre. Cofinality of the good partial colimits in every fibre then gives their cofinality in the total poset. The final diagram so obtained need not itself be continuous at all small directed suprema, exactly as Grothendieck observes on page~95. The reindexing on page~96 is therefore still needed.

The counterexample concerns this particular image-size formulation. The corrected small-subcategory construction is the familiar filtered-diagram argument, already used in SGA~4 I, 9.13.1. We did not find the specific counterexample above in the literature checked; it records our verification of the printed hypothesis, not a claim of priority.
\end{dcomment}


\msp{page 96}Mais revenons à ce cas
\[ X \;=\; \varinjlim X_i \;\left|\; \begin{array}{l} X_i\in M^\pi \\[1.5ex] I \text{ grand devant } \pi, \end{array}\right. \]
et introduisons à nouveau
\[ \tilde{I} \;=\; \{I'\subset I \mid I' \text{ filtrant, de cardinal } \le\pi\}, \]
donc cette fois $\tilde{I}$ est bien adapté à $\pi$ (filtrant, et toute partie filtrante de cardinal $\le\pi$ a une borne supérieure), et si on pose
\[ X_{\tilde{\imath}} \;=\; \varinjlim_{i\in I'_{\tilde{\imath}}} X_i, \]
on voit que le système inductif
\[ (X_{\tilde{\imath}})_{\tilde{\imath}\in\tilde{I}} \]
est parfaitement adapté à $\pi$ et la filtration cardinale. De plus, je dis que l'ensemble $\tilde{I}_M$ est \textit{cofinal} dans $\tilde{I}$. En effet, $\tilde{I}_M$ contient le sous-ensemble $\tilde{I}_0$ formé des $I'\in\tilde{I}$ qui ont un plus grand élément -- et $\tilde{I}_0$ est bien cofinal ! Maintenant, on peut appliquer
\printedpage{52}
% head: A. Grothendieck : Dérivateurs, Ch. XIII, Édition des Universités de Montpellier II et Paris VII
% foot: Version 25 janvier 2004
\msp{page 97}le ``il suffit'' dans \textit{b}), qui conclue que $X = \varinjlim X_i\in M$, q.e.d. % [sic] conclue -> conclut, verb form (indicative of conclure)

On a prouvé $b\Longrightarrow a$, sans utiliser l'hypothèse que tout objet de $\mathrm{Filt}^\pi$ est $\pi$-accessible. Mais on a utilisé cette hypothèse pour prouver que $b$ (ou $a$) implique que $M$ est stable par $\varinjlim$ grandes devant $\pi$ (lemme 2, p. 90).

\dgrblock{dgrb737}{The implication that does not need presentable objects}
\begin{dcomment}{exposition}{p97:On a prouvé $b\Longrightarrow a$}{The implication that does not need presentable objects}
The short argument on pages~96--97 is more robust than the intervening attempt to prove closure under large colimits. It proves $b_\pi\Rightarrow a_\pi$ directly, using only the strict filtration and the existence of perfectly adapted presentations. In particular it should not inherit the tacit presentability hypothesis needed for the preceding comma-category construction.

First suppose $X\in M$. Choose a perfectly adapted presentation $X=\varinjlim_I X_i$ at $\pi$. By $b_\pi$, the subset $I_M$ is cofinal. A cofinal subset of a $\pi^+$-directed poset is again $\pi^+$-directed: take an upper bound in the original poset for any at most $\pi$ elements, and then move above it into the cofinal subset. Restricting to $I_M$ therefore presents $X$ as a $\pi^+$-directed colimit of objects of $M^\pi$.

For the other implication in $a_\pi$, suppose $X=\varinjlim_{i\in I}Y_i$ with $Y_i\in M^\pi$ and $I$ grand devant $\pi$. Let $P$ be the poset of nonempty directed subsets $S\subset I$ of cardinal at most $\pi$, and put $Z_S=\varinjlim_S Y_i$. Strictness places each $Z_S$ in $\Filt^\pi$. The union construction makes $(Z_S)_{S\in P}$ perfectly adapted, with colimit $X$.

The key cofinal family consists of the subsets having a largest element. For any $S\in P$, choose an upper bound $s$ of all of $S$ in $I$ and put $S^+=S\cup\{s\}$. Then $S^+\in P$, it contains $S$, and
\[
 Z_{S^+}\simeq Y_s\in M^\pi.
\]
Hence $P_M$ is cofinal in $P$. The sufficiency direction in $b_\pi$ gives $X\in M$. This completes both directions of $a_\pi$.

The earlier inference on page~90 applied $b_\pi$ directly to an arbitrary directed presentation and omitted its continuity hypothesis; footnote~42 explicitly retracts that inference. The argument just given is the source's own repair. After it is established, the additional hypothesis that $M^\pi$ consists of $\pi^+$-presentable objects permits the closure argument through canonical comma categories. These are two different uses of the hypotheses, and retaining that separation prevents an artificial circularity in the proof.
\end{dcomment}


$a\overset{?}{\Longrightarrow} b$. Le ``il suffit'' dans \textit{b}) est conséquence immédiate de \textit{a}). Il reste à prouver que si, dans le système inductif $\pi$-adapté $(X_i)_{i\in I}$ de limite $X$, on a $X\in M$, alors $I_M$ est cofinal dans $I$. Mais par hypothèse \textit{a}), on aura $X = \varinjlim\limits_J Y_j$, avec $J$ grand devant $\pi$, les $Y_j$ dans $M^\pi$. Admettons que $\mathrm{Filt}^\pi\subset E_\pi$, \textit{i.e.} tout élément de $\mathrm{Filt}^\pi$, en particulier les $X_i$, sont $\pi$-accessibles.

Je pense que je n'arriverai pas à prouver $a_\pi\Longrightarrow b_\pi$, mais seulement ceci : Si $\pi'$ est un cardinal $>\pi$ (disons le successeur de $\pi$ dans la suite ordinale des cardinaux), alors $a_\pi\Longrightarrow b_\pi$.

\msp{page 98}Ceci va résulter du

\textbf{Lemme 3.} \textit{Soient $X\simeq\varinjlim\limits_I X_i$, $X\simeq\varinjlim\limits_J Y_j$ deux représentations de $X$ sous forme des $\varinjlim$ filtrantes, avec $I$ grand devant $\pi'$, $J$ grand devant $\pi$, avec $\pi$, $\pi'$ deux cardinaux tels que $\pi'>\pi$. On suppose $(X_i)_{i\in I}$ $\pi'$-adapté,} i.e. \textit{pour toute partie filtrante $I'$ de $I$ de cardinal $\le\pi'$, $I'$ a une borne supérieure $i_{I'}$, et $X_{i_{I'}} = \varinjlim\limits_{i\in I'} X_i$. Alors pour tout $i\in I$, on peut trouver une partie $I'$ de $I$, contenant $i$,} grande devant $\pi$, \textit{de cardinal $\le\pi'$, et une partie $J'$ de $J$, enfin un isomorphisme $I'\xrightarrow[\sim]{\lambda}J'$ et une application strictement croissante $\mu : J'\longrightarrow I'$ et des homomorphismes} \texttt{[fin du chapitre]}

\dgrblock{dgrb738}{Where the manuscript ends}
\begin{dcomment}{exposition}{p98:et une application strictement croissante}{Where the manuscript ends}
Grothendieck has reached the comparison of two presentations when the manuscript breaks off. The final words announce homomorphisms without specifying them. We therefore have a programme for a comparison lemma, not a completed statement whose omitted proof can simply be supplied. The separate larger-cardinal repair is a theorem with its own stated hypotheses; it does not turn this fragment retrospectively into a finished argument.

The material that follows belongs to the edition's apparatus. Its notation and terminology indices refer to manuscript page numbers. Thus the reference ``77'' beside cardinal filtrations points to manuscript page~77, on printed page~42 of this edition. The glossary then restates definitions from the whole chapter, including sections outside these pages. Its axiom list is particularly useful here: it records finite limits and colimits in \textbf{M~0}, two factorizations in \textbf{M~2}, the properness conditions in $\mathbf M''$, and the separate requirement $W=TF\circ TC$ in \textbf{M~2\,bis}. These distinctions explain why the names of the various Quillen structures cannot be replaced indiscriminately by the single modern phrase ``model category''.

The glossary is a guide to those definitions, and should be checked against their explicit formulation in the manuscript. In its account of a strict cardinal filtration, the qualification ``of cardinality at most $\pi$'' has disappeared; manuscript page~86, footnote~41, supplies it. Its final paragraph also says ``well adapted'' where the immediately preceding definition requires ``perfectly adapted''. The audit registers these two discrepancies separately, retaining the printed wording. The bound matters: by the representation axiom every object of $E$ is a filtered colimit of objects in any fixed stage, so unrestricted filtered-colimit closure would identify that stage with all of $E$. The printed glossary would then express a substantially different definition.

Finally, the bibliography after the glossary is the edition's bibliography. It includes the 1996 publication of Verdier's thesis; its presence is not evidence that Grothendieck had consulted that publication while writing these pages in 1990--1991. Our later references and our repairs are placed in the commentary and its own bibliography, where the date and the responsibility for each statement remain visible.
\end{dcomment}


\dgrblock{dgrb739}{A complete comparison after one successor step}
\begin{dcomment}{development}{p98:et des homomorphismes}{A complete comparison after one successor step}
The chapter stops while Grothendieck is constructing a comparison between two presentations. The final Lemma~3 has no completed conclusion or proof, and we do not supply words to its unfinished sentence. His preceding proposal does, however, suggest a precise statement which can be proved under explicit presentability assumptions. The printed implication on page~97 repeats $b_\pi$; the surrounding reference to a larger cardinal indicates $b_{\pi'}$ as the intended target.

\textbf{Successor comparison.}\ours
Let $E$ carry a strict cardinal filtration and let $\pi\geq\pi_0$ be infinite. Put $\kappa=\pi^+$ and $\mu=\kappa^+$. Assume that objects of $\Filt^\pi(E)$ are $\kappa$-presentable and that objects of $\Filt^\kappa(E)$ are $\mu$-presentable. Let $M$ be a replete full subcategory satisfying $a_\pi$. Then $M$ satisfies $b_\kappa$, that is, $b_{\pi^+}$ in the manuscript's notation.

We prove the statement rather than infer it from a general assertion that increasing a cardinal preserves every accessibility property. The literature checked includes SGA~4 I, 9.18, Grothendieck's later Chapter~XVIII, \S4.13 and Lemma~4.16.7.11, and \citet[\S8]{LiebermanRosickyVasey2020}. The argument below is a formulation for the particular continuous presentations used here; its use of alternating factorizations belongs to the standard comparison method for filtrations. We make no priority claim.

\emph{Proof.} The canonical-slice argument above shows first that $M$ is closed under $\kappa$-directed colimits. Indeed every object of $M$ has a presentation by objects of $M^\pi$, these objects are $\kappa$-presentable, and a $\kappa$-directed diagram of such presentations can be combined into one $\kappa$-filtered diagram of their small terms. The closure half of $a_\pi$ then applies.

We next obtain a more highly directed presentation of each $X\in M$, while keeping its terms in $M$. Write $X=\varinjlim_{j\in J}Y_j$ as in $a_\pi$. Any subset $S\subset J$ of cardinal at most $\kappa$ lies in a $\kappa$-directed subposet $K\subset J$ of size at most $\kappa$. To see this, first adjoin any element of the nonempty poset $J$ if $S$ is empty, then enumerate $S$ in length at most $\kappa$, repeating elements if necessary. At stage $\alpha<\kappa$, choose an upper bound in $J$ for the current enumerated element and all previously chosen bounds. There are at most $\pi$ of them. The increasing sequence of bounds, together with $S$, is $\kappa$-directed: any fewer than $\kappa$ of its elements is accounted for before some stage, since $\kappa$ is regular.

Let $H$ be the poset of all these $\kappa$-directed subposets of $J$ of size at most $\kappa$. It is $\mu$-directed. The union of at most $\kappa$ of its members has size at most $\kappa$, and the construction just given places that union in another member of $H$. Put
\[
 Z_K=\varinjlim_{j\in K}Y_j\qquad(K\in H).
\]
The closure already proved gives $Z_K\in M$, and strictness of the filtration gives $Z_K\in\Filt^\kappa(E)$. Thus $Z_K$ is $\mu$-presentable by hypothesis. The union-of-subdiagrams argument gives $X\simeq\varinjlim_{K\in H}Z_K$.

Now take a perfectly adapted presentation $X=\varinjlim_{i\in I}X_i$ at cardinal $\kappa$. Its index $I$ is $\mu$-directed and its objects are $\mu$-presentable. Fix $i_0\in I$. We will construct increasing sequences $(i_\alpha)_{\alpha<\kappa}$ in $I$ and $(K_\alpha)_{\alpha<\kappa}$ in $H$, together with maps over $X$
\[
 X_{i_\alpha}\xrightarrow{u_\alpha}Z_{K_\alpha}
       \xrightarrow{v_\alpha}X_{i_{\alpha+1}},
\]
such that
\[
 v_\alpha u_\alpha=X(i_\alpha\leq i_{\alpha+1}),
\]
and, for $\alpha<\beta<\kappa$,
\[
 u_\beta X(i_{\alpha+1}\leq i_\beta)v_\alpha
       =Z(K_\alpha\subset K_\beta).                 \tag{SC}
\]
The first equality makes a trip through the other presentation agree with a transition in the $X$-diagram; the second does the same for the $Z$-diagram. They will give inverse maps on colimits.

Here are all the inductive choices. At a limit stage $\beta<\kappa$, choose $i_\beta$ above the earlier $i_\alpha$ and $i_{\alpha+1}$; at a successor it has already been chosen. Since $X_{i_\beta}$ is $\mu$-presentable, its map to $X=\varinjlim_H Z_K$ factors through some $Z_K$. Enlarge $K$ above all the previously chosen $K_\alpha$. For each $\alpha<\beta$, the two maps from $Z_{K_\alpha}$ to this stage occurring in (SC) agree after passage to $X$. Presentability of $Z_{K_\alpha}$ makes that equality hold at some later $H$-stage. There are fewer than $\kappa$, hence fewer than $\mu$, such equalities. The $\mu$-directedness of $H$ permits one enlargement making all of them hold. Call the resulting index $K_\beta$ and the factorization $u_\beta$.

Next factor $Z_{K_\beta}\to X$ through an $X$-stage, using its $\mu$-presentability. Enlarge that stage above $i_\beta$. The composite with $u_\beta$ and the transition out of $X_{i_\beta}$ have equal images in $X$; presentability of $X_{i_\beta}$ allows a further enlargement on which they are equal. This gives $i_{\beta+1}$ and $v_\beta$. The same process starts at $\beta=0$, when there are no equations (SC) from earlier stages. Every choice takes place in a small diagram, and the regularity and directedness just used justify every limit stage.

Equations (SC) imply that the $u_\alpha$ define a map between the two chain colimits: compose (SC) with $u_\alpha$ and use $v_\alpha u_\alpha=X(i_\alpha\leq i_{\alpha+1})$. They also imply that the $v_\alpha$, followed by the structural maps to the $X$-chain colimit, are compatible. The two resulting maps are inverse, by the same displayed equalities. Consequently
\[
 \varinjlim_{\alpha<\kappa}X_{i_\alpha}
 \simeq\varinjlim_{\alpha<\kappa}Z_{K_\alpha}\in M.
\]
The membership assertion uses closure of $M$ under $\kappa$-directed colimits, since the regular ordinal $\kappa$ is a $\kappa$-directed poset.

The perfectly adapted $X$-diagram has a supremum $i_* =\sup_{\alpha<\kappa}i_\alpha$ and is continuous there: the image of this chain is a directed subset of $I$ with at most $\kappa$ elements. Therefore $X_{i_*}\simeq\varinjlim_{\alpha<\kappa}X_{i_\alpha}$ belongs to $M$, and $i_*\geq i_0$. Thus $I_M$ is cofinal whenever $X\in M$. Conversely, if $I_M$ is cofinal then it is $\mu$-directed, hence $\kappa$-directed, and closure of $M$ gives $X=\varinjlim_{I_M}X_i\in M$. This proves $b_\kappa$.\hfill$\square$

The cardinal shift has a definite function. It permits a comparison chain of length $\pi^+$ whose final colimit remains a single stage of the larger adapted presentation. At the original cardinal, that chain is too large for the continuity clause. The counterexample with successor and limit ordinals shows that this obstruction can be real. No continuum hypothesis entered the proof, but presentability at both filtration levels did; the filtration axioms alone do not provide it.
\end{dcomment}


\dgrblock{dgrb740}{Canonical slices, later accessibility, and weak equivalences}
\begin{dcomment}{later}{p98:et des homomorphismes}{Canonical slices, later accessibility, and weak equivalences}
Grothendieck returned to these questions at much greater length in Chapter~XVIII. Definition~4.13.1 there asks that a full subcategory be replete, accessible, and accessibly included in the ambient accessible category. Lemma~4.16.7.11 tests membership by cofinality of the small canonical slice inside the ambient small-object slice. That is an intrinsic form of the generation condition, with presentability built into the indexing category. It does not assert that every continuous presentation already has cofinally many stages in the subcategory at the original rank. The two powerset examples above therefore leave the later criterion intact. The edition consulted is dated 17~December~2025; the cited lemma lies on manuscript pages~488--490 \citep{Grothendieck1990XVIII}.

In present terminology, accessibility and accessible inclusion are distinct conditions, and a rank belongs to the statement. The definition and the description as a completion under directed colimits are given by \citet[Definition~2.1 and Theorem~2.26]{AdamekRosicky1994}. The earlier monograph of \citet{MakkaiPare1989} situates these constructions within categorical model theory. For the higher-categorical formulation, \citet[Proposition~5.4.2.2]{Lurie2009HTT} characterizes accessible $\infty$-categories by small generating subcategories of $\kappa$-compact objects for a suitable regular $\kappa$. Its nearby rank-enlargement results retain explicit cardinal conditions; a bare change from one regular cardinal to a larger one is not a general theorem.

The modern theory of filtrations also distinguishes a directed presentation from a continuous chain. \citet[\S8, especially Lemma~8.4 and the following paragraph]{LiebermanRosickyVasey2020} discusses continuous filtrations and their agreement on a closed unbounded set of indices when the size has the appropriate regularity. This helps explain the example above: the successor-ordinal presentation of $\omega_1$ is not continuous at its countable limit stages. The comparison theorem for two smooth filtrations therefore cannot turn its stages into stages of the continuous limit-ordinal presentation. We checked that discussion and Theorem~8.8 in the authors' 5~June~2019 version; we do not identify its notion of filtration with Grothendieck's cardinal filtration of the whole category.

For homotopical algebra, the relevant application is to a full subcategory of the arrow category. \citet[Proposition~2.2 in author version~2; Proposition~2.3 in the journal]{RaptisRosicky2015} gives a controlled accessibility rank for weak equivalences in a combinatorial model category after the replacement functors preserve the specified presentable objects. Their Theorem~4.2 supplies better bounds under additional hypotheses, and Corollary~5.1 proves finite accessibility for weak equivalences of simplicial sets. These are genuine refinements of the small-generation problem motivating this chapter. They do not show that weak equivalences in every combinatorial model category are finitely accessible, nor do they establish the source's universal-presentation condition for an arbitrary filtration. The present status of the local questions is thus precise: $b_\pi\Rightarrow a_\pi$ is proved in the manuscript; the three-way equivalence as stated is false; the unfinished comparison has the explicit conditional completion proved above.
\end{dcomment}


\newpage
\printedpage{53}
% head: A. Grothendieck : Dérivateurs, Ch. XIII, Édition des Universités de Montpellier II et Paris VII
% foot: Version 25 janvier 2004

\section*{Index des notations}

\idx{LLP, RLP}{1}
\idx{$\Phi_*$, $\Psi^*$}{1}
\idx{$\widetilde{\Phi}$, $\overline{\Psi}$}{1}
\idx{$\Phi \longleftrightarrow \Psi$}{3}
\idx{$\mathrm{Cof}_W$, $\mathrm{Fib}_W$}{8}
\idx{$W^{\mathrm{univ}}$}{9}
\idx{$W_{\mathrm{univ}}$}{10}
\idx{$\underline{\mathrm{Fib}}_W S\subset\mathcal{M}/S$}{10}
\idx{$W_S$}{10}
\idx{$W_S^{\mathrm{u}}$}{10}
\idx{$S\backslash\underline{\mathrm{Cof}}_W$}{10}
\idx{$\Phi \barlongleftrightarrow \Psi$}{28}
\idx{$\mathcal{M}_{\mathrm{cf}}$, $W_{\mathrm{cf}}$}{29}
\idx{$\mathrm{Ho}_W$}{29}
\idx{$\mathrm{Hom}_W$}{30}
\idx{$\pi_{\mathrm{d}}$, $\pi$}{30}
\idx{$\mathcal{M}_{\mathrm{c}}$}{30}
\idx{$W_{\mathrm{c}}$}{31}
\idx{$\pi_{\mathrm{g}}$}{31}
\idx{$\mathcal{M}_{\mathrm{f}}$}{31}
\idx{$f\sim g$}{32}
\idx{$\Psi,\ \Phi,\ \mathcal{M}(\Psi),\ W(\Psi),\ \mathcal{M}(\Phi),\ W(\Phi)$}{33}
\idx{$\mathcal{M}(I)$, $W(I)$}{34}
\idx{$F(I),\ TF(I),\ C(I),\ TC(I)$}{35}
\idx{$\alpha_s$}{35}
\idx{$u^*$, $u_*$, $u_!$}{45}
\idx{$u^*_W$, $u^W_*$, $u^W_!$}{49}
\idx{$\mathrm{HO}_W(I)$}{49}
\idx{Diag}{50}
\idx{$\mathrm{HO}_{\mathcal{M},W}(I)$, $\mathbf{D}_{\mathcal{M},W}(I)$}{54}
\idx{$\mathrm{HO}^{\mathrm{lc}}_W(I)$, $u^*_{\mathrm{lc}}$}{56}
\idx{$u^{\mathrm{f}}_*$}{59}
\idx{$u^*_{\mathrm{c}}$}{67}
\idx{$\mathfrak{U}$}{77}
\idx{$\mathrm{Filt}^\pi(E)$}{77}

% ---- from Der-ChXIII-25-01-04_p056-060.tex
\printedpage{54}
% head: A. Grothendieck : Dérivateurs, Ch. XIII, Édition des Universités de Montpellier II et Paris VII
% foot: Version 25 janvier 2004

\entry{$L_\pi$}{77}
\entry{$\tilde{I}$}{79}
\entry{$\underline{M}$}{86}
\entry{$I_M$}{87}
\entry{$M^\pi$}{87}

\clearpage

\printedpage{55}
% head: A. Grothendieck : Dérivateurs, Ch. XIII, Édition des Universités de Montpellier II et Paris VII
% foot: Version 25 janvier 2004

\section*{Index terminologique}

\entry{accessible (partie close à gauche)}{75}
\entry{admissible (à gauche)}{59}
\entry{catégorie de modèles}{7}
\entry{catégorie de Quillen close}{25}
\entry{clos à gauche, à droite}{2}
\entry{couple de Quillen}{5}
\entry{faiblement saturée}{7}
\entry{filtration cardinale}{77}
\entry{filtration cardinale stricte}{86}
\entry{flèche coüniversellement dans $W$}{9}
\entry{flèche universellement dans $W$}{10}
\entry{grand devant $\pi$}{77}
\entry{objet chemin}{32}
\entry{objet cylindre}{32}
\entry{$\pi$-accessible (partie d'une catégorie munie d'une filtration cardinale stricte)}{86}
\entry{pré-Q-catégorie de modèles à gauche, à droite}{12}
\entry{pré-Quillen à gauche, à droite}{14}
\entry{pré-quintuplet de Quillen clos}{21}
\entry{propre}{8}
\entry{propriété de factorisation de Quillen}{5}
\entry{Q-catégorie de modèles}{8}
\entry{Q-catégorie de modèles à droite}{21}
\entry{Q-catégorie de modèles “close”}{16}
\entry{Q-homotopie à gauche, à droite}{30}
\entry{quillénisable}{50}
\entry{quillénisable sur Diag}{51}
\entry{quillénisable sur Diag uniformément et exactement à gauche}{57}
\entry{quillénisable uniformément à gauche sur Diag}{57}
\entry{quillénisation uniforme à droite sur Diag}{58}
\entry{quillénisation uniforme à gauche sur Diag}{57}
\entry{quillénisation uniforme et exacte à gauche sur Diag}{57}
\entry{représentation bien adaptée}{84}
\entry{représentation parfaitement adaptée}{86}
\entry{strictement admissible}{61}
\entry{structure de Quillen}{28}

\clearpage

\printedpage{56}
% head: A. Grothendieck : Dérivateurs, Ch. XIII, Édition des Universités de Montpellier II et Paris VII
% foot: Version 25 janvier 2004

\entry{structure de Quillen exacte}{27}
\entry{structure de Quillen exacte à gauche, à droite}{27}
\entry{structure de Quillen faible}{27}
\entry{structure de Quillen faible biexacte}{28}
\entry{structure de Quillen faible close, close à gauche, à droite}{28}
\entry{structure de Quillen parfaite}{28}
\entry{système de Quillen clos}{26}
\entry{système de Quillen faible}{26}
\entry{système de Quillen faible exact à gauche, à droite}{26}
\entry{système inductif bien adapté}{84}
\entry{système inductif parfaitement adapté}{86}
\entry{triple clos de Quillen}{25}
\entry{triple de Quillen}{14}
\entry{$W$-cofibration}{8}
\entry{$W$-équivalence coüniverselle}{9}
\entry{$W$-fibration}{8}

\clearpage

\printedpage{57}
% head: A. Grothendieck : Dérivateurs, Ch. XIII, Édition des Universités de Montpellier II et Paris VII
% foot: Version 25 janvier 2004

\section*{Liste des axiomes}

\noindent $\mathcal{M}$ catégorie

\noindent $W \ \subset\ \Fl(\mathcal{M})$ , \qquad $C$, $F$, $TC$, $TF \ \subset\ \Fl(\mathcal{M})$

\begin{description}[leftmargin=4.5em,labelwidth=4em]
\item[M 0] La catégorie $\mathcal{M}$ admet des limites projectives finies et des limites inductives finies.

\item[M 1] Les flèches appartenant à $C$ ont la propriété de relèvement à gauche relativement à celles appartenant à $TF$, et les flèches appartenant à $F$ ont la propriété de relèvement à droite relativement à celles appartenant à $TC$.

\item[M 2] Toute flèche de $\mathcal{M}$ se décompose en une flèche de $C$ suivie d'une flèche de $TF$, et en une flèche de $TC$ suivie d'une flèche de $F$.

\item[M 3] La partie $F$ de $\Fl(\mathcal{M})$ est stable par composition et changement de base, $C$ est stable par composition et cochangement de base, $F$ et $C$ contiennent les isomorphismes.

\item[M 4] La partie $TF$ de $\Fl(\mathcal{M})$ est stable par composition et changement de base, $TC$ est stable par composition et cochangement de base, $TF$ et $TC$ contiennent les isomorphismes.

\item[M 5] Si dans un triangle commutatif de $\mathcal{M}$ deux des trois flèches sont dans $W$, il en est de même de la troisième ; $W$ contient les isomorphismes.

\item[M$'$] $TC \ \subset\ C \cap W$, $TF \ \subset\ F \cap W$.

\item[M$''$] Les flèches appartenant à $C$ sont des $W$-cofibrations, et les flèches appartenant à $F$ sont des $W$-fibrations.

\item[M$'_{\mathbf{0}}$(g)] $TC = C \cap W$.

\item[M$'_{\mathbf{0}}$(d)] $TF = F \cap W$.

\item[M 2\,bis] $W = TF \circ TC$.
\end{description}

\clearpage

\printedpage{58}
% head: A. Grothendieck : Dérivateurs, Ch. XIII, Édition des Universités de Montpellier II et Paris VII
% foot: Version 25 janvier 2004

\setcounter{footnote}{0}
\renewcommand{\thefootnote}{\fnsymbol{footnote}}
\section*{Glossaire}

\subsection*{Section 1}

\begin{itemize}
\item $\Phi \ \subset\ \Fl(\mathcal{M})$

$\Phi$ est \textit{clos à gauche} : il existe \ $\Psi \ \subset\ \Fl(\mathcal{M})$ tel que $\Phi$ soit formé des flèches ayant la propriété de relèvement à gauche relativement aux flèches de $\Psi$.

\item $\Psi \ \subset\ \Fl(\mathcal{M})$

$\Psi$ est \textit{clos à droite} : il existe \ $\Phi \ \subset\ \Fl(\mathcal{M})$ tel que $\Psi$ soit formé des flèches ayant la propriété de relèvement à droite relativement aux flèches de $\Phi$.

\item $\Phi$, $\Psi \ \subset\ \Fl(\mathcal{M})$

$(\Phi, \Psi)$ a la \textit{propriété de factorisation de Quillen} :
\begin{itemize}
\item[a)] les flèches de $\Phi$ ont la propriété de relèvement à gauche relativement aux flèches de $\Psi$ ;
\item[b)] toute flèche $f$ de $\mathcal{M}$ se factorise en
\[
f = pi \quad , \qquad i \in \Phi \ , \quad p \in \Psi \ .
\]
\end{itemize}

\item $(\Phi, \Psi)$ \textit{couple de Quillen} :
\begin{itemize}
\item[a)] $(\Phi, \Psi)$ a la propriété de factorisation de Quillen ;
\item[b)] $\Phi$ est clos à gauche et $\Psi$ clos à droite.
\end{itemize}
\end{itemize}

\subsection*{Section 2}

\begin{itemize}
\item $W \ \subset\ \Fl(\mathcal{M})$

\textit{$W$ faiblement saturé} :
\begin{itemize}
\item[a)] $W$ contient les isomorphismes ;
\item[b)] si deux parmi les trois flèches d'un triangle commutatif sont dans $W$, la troisième aussi.
\end{itemize}

\item $W \ \subset\ \Fl(\mathcal{M})$

Un morphisme $p : X \longrightarrow Y$ de $\mathcal{M}$ est une \textit{$W$-fibration} ($p \in \mathrm{Fib}_W$) si pour tout diagramme de carrés cartésiens
\end{itemize}

% ---- from Der-ChXIII-25-01-04_p061-065.tex
\printedpage{59}
% head: A. Grothendieck : Dérivateurs, Ch. XIII, Édition des Universités de Montpellier II et Paris VII

\[
\begin{tikzcd}
X \arrow[d, "p"'] & X' \arrow[l] \arrow[d] & X'' \arrow[l, "t"'] \arrow[d] \\
Y & Y' \arrow[l] & Y'' \arrow[l, "s"]
\end{tikzcd}
\quad ,
\]

si $s \in W$, alors $t \in W$.

\begin{itemize}
\item[\textbullet] $W \ \subset\ \Fl(\Mo)$

Un morphisme $i : A \longrightarrow B$ de $\Mo$ est une $W$-\textit{cofibration} ($i \in \mathrm{Cof}_W$) si pour tout diagramme de carrés cocartésiens
\[
\begin{tikzcd}
A'' \arrow[d] & A' \arrow[l, "s"'] \arrow[d] & A \arrow[l] \arrow[d, "i"] \\
B'' & B' \arrow[l, "t"] & B \arrow[l]
\end{tikzcd}
\quad ,
\]
si $s \in W$, alors $t \in W$.

\item[\textbullet] $W \ \subset\ \Fl(\Mo)$

Un morphisme $p : X \longrightarrow Y$ de $\Mo$ est \textit{universellement dans} $W$, ou une $W$-\textit{équivalence universelle} ($p \in W_{\mathrm{univ}}$) si pour tout carré cartésien
\[
\begin{tikzcd}
X \arrow[d, "p"'] & X' \arrow[l] \arrow[d, "p'"] \\
Y & Y' \arrow[l]
\end{tikzcd}
\quad ,
\]
la flèche $p'$ appartient à $W$.

\item[\textbullet] $W \ \subset\ \Fl(\Mo)$

Un morphisme $i : A \longrightarrow B$ de $\Mo$ est \textit{coüniversellement dans} $W$, ou une $W$-\textit{équivalence coüniverselle} ($i \in W^{\mathrm{univ}}$) si pour tout carré cocartésien
\[
\begin{tikzcd}
A' \arrow[d, "i'"'] & A \arrow[l] \arrow[d, "i"] \\
B' & B \arrow[l]
\end{tikzcd}
\quad ,
\]
la flèche $i'$ appartient à $W$.
\end{itemize}
% foot: Version 25 janvier 2004

\printedpage{60}
% head: A. Grothendieck : Dérivateurs, Ch. XIII, Édition des Universités de Montpellier II et Paris VII

\begin{itemize}
\item[\textbullet] \textit{Pré-Q-catégorie de modèles à gauche} (resp. \textit{à droite}) :

$\Mo$ et $W$, $C$, $F$, $TC$, $TF$ $\subset$ $\Fl(\Mo)$ satisfaisant aux axiomes \textbf{M~0} à \textbf{M~5}, $\mathbf{M}'$, $\mathbf{M}''$ et \Mzg{} (resp. \Mzd).

\item[\textbullet] \textit{Q-catégorie de modèles}, \textit{triple de Quillen} :

catégorie de modèles au sens de Quillen.

\item[\textbullet] \textit{Q-catégorie de modèles “close”}, \textit{triple clos de Quillen}, \textit{catégorie de Quillen close}, \textit{système de Quillen clos} :

catégorie de modèles fermée au sens de Quillen.

\item[\textbullet] \textit{Pré-Quillen à gauche} (resp. \textit{à droite}) :

pré-Q-catégorie de modèles à gauche (resp. à droite).

\item[\textbullet] \textit{Pré-quintuplet de Quillen clos} :

$W$, $C$, $F$, $TC$, $TF$ $\subset$ $\Fl(\Mo)$ satisfaisant aux axiomes \textbf{M~0} à \textbf{M~5}, $\mathbf{M}'$, $\mathbf{M}''$, $C$, $TC$ étant clos à gauche et $F$, $TF$ clos à droite.

\item[\textbullet] \textit{Système de Quillen faible} :

$W$, $C$, $F$, $TC$, $TF$ $\subset$ $\Fl(\Mo)$ satisfaisant aux axiomes \textbf{M~0} à \textbf{M~5}, $\mathbf{M}'$, $\mathbf{M}''$.

\item[\textbullet] \textit{Système de Quillen faible exact à gauche} (resp. \textit{à droite}) :

$W$, $C$, $F$, $TC$, $TF$ $\subset$ $\Fl(\Mo)$ satisfaisant aux axiomes \textbf{M~0} à \textbf{M~5}, $\mathbf{M}'$, $\mathbf{M}''$ et \Mzg{} (resp. \Mzd).
\end{itemize}

\section*{Section 3}

\begin{itemize}
\item[\textbullet] $(\Mo,W)$ \textit{catégorie de modèles} :

$\Mo$ catégorie satisfaisant à l'axiome \textbf{M~0} ;

$W$ $\subset$ $\Fl(\Mo)$ faiblement saturé.

\item[\textbullet] \textit{Structure de Quillen faible associée à la catégorie de modèles} $(\Mo,W)$ :

quadruplet $(C,F,TC,TF)$ de parties de $\Fl(\Mo)$ satisfaisant aux axiomes \textbf{M~1} à \textbf{M~5}, $\mathbf{M}'$, $\mathbf{M}''$.

\item[\textbullet] \textit{Structure de Quillen faible exacte} :

structure de Quillen faible satisfaisant en outre à l'axiome \textbf{M~2\,bis}.

\item[\textbullet] \textit{Structure de Quillen faible exacte à gauche} (resp. \textit{à droite}) :

structure de Quillen faible satisfaisant en outre à l'axiome \Mzg{} (resp. \Mzd).
\end{itemize}
% foot: Version 25 janvier 2004

\printedpage{61}
% head: A. Grothendieck : Dérivateurs, Ch. XIII, Édition des Universités de Montpellier II et Paris VII

\begin{itemize}
\item[\textbullet] \textit{Structure de Quillen faible close à gauche} (resp. \textit{à droite}) :

structure de Quillen faible $(C,F,TC,TF)$ telle que $C$, $TC$ (resp. $F$, $TF$) soient clos à gauche (resp. à droite).

\item[\textbullet] \textit{Structure de Quillen faible close} :

structure de Quillen faible à la fois close à gauche et à droite.

\item[\textbullet] \textit{Structure de Quillen} :

structure de Quillen faible “biexacte” :

structure de Quillen faible à la fois exacte à gauche et à droite

$=$ catégorie de modèles de Quillen propre.

\item[\textbullet] \textit{Structure de Quillen parfaite} :

structure de Quillen faible “biexacte” et close

$=$ catégorie de modèles fermée de Quillen propre.
\end{itemize}

\section*{Section 5}

\begin{itemize}
\item[\textbullet] Catégorie de modèles $(\Mo,W)$ \textit{quillénisable} :

il existe une structure de Quillen faible \textit{exacte} $(C,F,TC,TF)$ associée à $(\Mo,W)$.

\item[\textbullet] $(\Mo,W)$ \textit{quillénisable sur} Diag $\subset$ Cat :

pour tout $I$ dans Diag, la catégorie des modèles $(\Mo(I),\, W(I))$ est quillénisable.
\end{itemize}

\section*{Section 6}

\begin{itemize}
\item[\textbullet] Catégorie de modèles $(\Mo,W)$ \textit{quillénisable uniformément à gauche sur} Diag :

il existe une structure de Quillen faible \textit{close} $(C,F,TC,TF)$ pour $(\Mo,\, W)$ telle que pour tout $I$ dans Diag, si l'on définit $C(I)$ et $TC(I)$ “argument par argument” (\footnote{Un morphisme de foncteurs $u : X \longrightarrow Y$ de $I^{\circ}$ vers $\Mo$ étant dans $C(I)$ (resp. $TC(I)$) si et seulement si pour tout objet $i$ de $I$, $u_i$ est dans $C$ (resp. $TC$).}), et $F(I)$ et $TF(I)$ par la propriété de relèvement à droite relativement à $TC(I)$ et $C(I)$ respectivement (\footnote{Un morphisme de $\Mo(I)$ étant dans $F(I)$ (resp. $TF(I)$) si et seulement si il a la propriété de relèvement à droite relativement aux flèches appartenant à $TC(I)$ (resp. $C(I)$).}), le quadruplet $(C(I),TC(I),F(I),TF(I))$ soit une structure de Quillen faible \textit{close exacte} pour $(\Mo(I),\, W(I))$. On dit alors que $(C,F,TC,TF)$ est une \textit{quillénisation uniforme à gauche} de $(\Mo,\, W)$ sur Diag).
\end{itemize}
% foot: Version 25 janvier 2004

\printedpage{62}
% head: A. Grothendieck : Dérivateurs, Ch. XIII, Édition des Universités de Montpellier II et Paris VII

\begin{itemize}
\item[\textbullet] Catégorie de modèles $(\Mo,W)$ \textit{quillénisable sur} Diag \textit{uniformément et exactement à gauche} :

$(\Mo,W)$ admet une quillénisation uniforme à gauche sur Diag $(C,F,TC,TF)$ telle que $TC = C \cap W$ (autrement dit, telle que $(C,F,TC,TF)$ soit une strucure de Quillen faible \textit{exacte à gauche}). On dit alors que $(C,F,TC,TF)$ est une \textit{quillénisation uniforme et exacte à gauche} de $(\Mo,\, W)$ sur Diag. % [sic] strucure -> structure, letter t omitted
\end{itemize}

\section*{Section 7}

\begin{itemize}
\item[\textbullet] $I$ ensemble ordonné filtrant, \quad $\pi$ cardinal

$I$ \textit{grand devant} $\pi$ :

toute partie de $I$ de cardinal inférieur ou égal à $\pi$ admet un majorant dans $I$.

\item[\textbullet] $E$ catégorie

Une \textit{filtration cardinale} de $E$ est une filtration croissante
\[
\Big(\Filt^{\pi}(E)\Big)_{\pi\geq\pi_0} \qquad , \qquad\qquad \pi_0 \text{ un cardinal infini fixé,}
\]
de $E$ par des sous-catégories strictement pleines $\Filt^{\pi}(E)$, indexée par les cardinaux $\pi \geq \pi_0$ appartenant à un univers donné, et satisfaisant aux conditions suivantes :
\begin{itemize}
\item[\textit{a})] pour tout $\pi \geq \pi_0$, $\Filt^{\pi}(E)$ est équivalente à une petite catégorie ;
\item[\textit{b})] la catégorie $E$ admet des limites inductives filtrantes indexées par des ensembles ordonnés $I$ grands devant $\pi_0$, et pour tout $\pi \geq \pi_0$, $\Filt^{\pi}(E)$ est stable par les dites limites si $\mathrm{card}\, I \leq \pi$ ;
\item[\textit{c})] Pour tout $\pi \geq \pi_0$, et pour tout $X \in \mathrm{Ob}\, E$, il existe un isomorphisme
\[
X \ \simeq\ \varinjlim_{i\in I} X_i \ ,
\]
où $(X_i)_{i\in I}$ est un système inductif filtrant dans $E$, indexé par un ensemble ordonné $I$ grand devant $\pi$, et où les $X_i$ sont des objets de $\Filt^{\pi}(E)$. De plus, si $X \in \mathrm{Ob}\,\Filt^{\pi'}(E)$, $\pi' \geq \pi$, on peut prendre $I$ tel que $\mathrm{card}\, I \leq \pi'^{\pi}$.
\end{itemize}

\item[\textbullet] $E$ catégorie, \quad $(\Filt^{\pi}(E))_{\pi\geq\pi_0}$ filtration cardinale de $E$, \quad $\pi \geq \pi_0$ cardinal

Un \textit{système inductif bien adapté} (à la filtration cardinale donnée, pour les cardinaux $\pi_0$, $\pi$ donnés) est un système inductif filtrant $(X_i)_{i\in I}$ satisfaisant aux conditions suivantes :
\begin{itemize}
\item[\textit{a})] $I$ est grand devant $\pi$, et pour tout $i \in I$, $X_i \in \Filt^{\pi}(E)$ ;
\item[\textit{b})] Toute partie $J$ de $I$ de cardinal $\leq \pi$, et grande devant $\pi_0$, admet une borne supérieure $i_J$, et on a
\[
X_{i_J} \ \xleftarrow{\ \sim\ }\ \varinjlim_{i\in J} X_i \ .
\]
% foot: Version 25 janvier 2004
\end{itemize}

\printedpage{63}
% head: A. Grothendieck : Dérivateurs, Ch. XIII, Édition des Universités de Montpellier II et Paris VII

Si $X$ est un objet de $E$, une \textit{représentation bien adaptée} de $X$, en termes d'objets de $\Filt^{\pi}(E)$, est la donnée d'un système inductif bien adapté $(X_i)_{i\in I}$, à la filtration cardinale donnée, pour les cardinaux $\pi_0$, $\pi$, et d'un isomorphisme
\[
X \ \xleftarrow{\ \sim\ }\ \varinjlim_{i\in I} X_i \ .
\]

\item[\textbullet] $E$ catégorie admettant des petites limites inductives filtrantes

Une \textit{filtration cardinale stricte} de $E$ est une filtration cardinale $(\Filt^{\pi}(E))_{\pi\geq\pi_0}$ telle que pour tout cardinal $\pi \geq \pi_0$, $\Filt^{\pi}(E)$ soit stable par limites inductives filtrantes.

\item[\textbullet] $E$ catégorie admettant des petites limites inductives filtrantes

$(\Filt^{\pi}(E))_{\pi\geq\pi_0}$ filtration cardinale stricte de $E$, \quad $\pi \geq \pi_0$ cardinal

Un \textit{système inductif parfaitement adapté} (à la filtration cardinale donnée, pour le cardinal $\pi$) est un système inductif filtrant $(X_i)_{i\in I}$ satisfaisant aux conditions suivantes :
\begin{itemize}
\item[\textit{a})] $I$ est grand devant $\pi$, et pour tout $i \in I$, $X_i \in \Filt^{\pi}(E)$ ;
\item[\textit{b})] Toute partie $J$ de $I$ filtrante de cardinal $\leq \pi$, admet une borne supérieure $i_J$, et on a
\[
X_{i_J} \ \xleftarrow{\ \sim\ }\ \varinjlim_{i\in J} X_i \ .
\]
\end{itemize}
Si $X$ est un objet de $E$, une \textit{représentation parfaitement adaptée} de $X$, en termes d'objets de $\Filt^{\pi}(E)$, est la donnée d'un système inductif bien adapté $(X_i)_{i\in I}$, à la filtration cardinale donnée, pour le cardinal $\pi$, et d'un isomorphisme
\[
X \ \xleftarrow{\ \sim\ }\ \varinjlim_{i\in I} X_i \ .
\]
\end{itemize}
% foot: Version 25 janvier 2004

% ---- from Der-ChXIII-25-01-04_p066-066.tex
\printedpage{64}
% head: A. Grothendieck : Dérivateurs, Ch. XIII, Édition des Universités de Montpellier II et Paris VII
% foot: Version 25 janvier 2004

\section*{Références}

\noindent H. J. Baues, ``Algebraic Homotopy'', Cambridge studies in advanced mathematics, vol. 15 (1989).

\noindent R. Hartshorne, ``Residues and Duality'', Lecture Notes in Mathematics, vol. 20, Springer-Verlag (1966).

\noindent D. G. Quillen, ``Homotopical Algebra'', Lecture Notes in Mathematics, vol. 43, Springer-Verlag (1967).

\noindent R. W. Thomason, \textit{Cat as a Closed Model Category}, Cahiers de Topologie et Géométrie Différentielle, Vol. XXI-3, pp. 305-324 (1980).

\noindent J.-L. Verdier, ``Des catégories dérivées des catégories abéliennes'', Astérisque, vol. 239 (1996).

\noindent SGA 4, I, ``Théorie des Topos et Cohomologie Étale des Schémas'', Séminaire de Géométrie Algébrique du Bois-Marie, dirigé par M. Artin, A. Grothendieck, J.-L. Verdier, Lecture Notes in Mathematics, vol. 269, Springer-Verlag (1972).
\begin{otherlanguage}{english}
\dgroverview{XIII-b}{Chapter XIII \textperiodcentered\ \textup{XIII-b} \textperiodcentered\ Overview}
\selectlanguage{english}

The second part of Chapter XIII joins two constructions that Grothendieck needs for the same purpose. He wants a homotopy theory to provide a coherent calculus of diagrams and their direct images. He also wants the replacements used in that calculus to be obtainable from a controlled set of small data. Sections~5 and~6 address the first problem through quillenisation; section~7 returns to the cardinal filtrations of SGA~4 to address the second. The manuscript ends during the statement of its final comparison lemma. The break is part of the surviving source, and the commentary does not supply a purported continuation of Grothendieck's text \citep{GrothendieckDerivateursXIII}.

Throughout these pages, $\mathcal M(I)$ means the category of contravariant diagrams $I^{\mathrm{op}}\to\mathcal M$, and $W(I)$ consists of objectwise weak equivalences. This variance explains both the right Kan extension formula over $I'/s$ and the fact that a greatest element of a poset computes the ordinary limit by evaluation. Grothendieck's phrase ``cat\'egorie de mod\`eles'' initially specifies much less than a modern Quillen model category: it is a category with finite limits and colimits and a two-out-of-three class of morphisms. Cofibrations, fibrations, closedness, and exactness are additional data or properties. Keeping those qualifications visible is necessary both to understand the proofs and to avoid attributing modern theorems to a weaker hypothesis.

A closed exact structure has more force than the subsequent doubts suggest. Exactness says that every weak equivalence factors as a trivial cofibration followed by a trivial fibration. Closedness supplies closure under retracts for the lifting classes. The elementary retract argument therefore gives $TC=C\cap W$ and $TF=F\cap W$. In a uniformly left quillenised diagram theory, restriction preserves the objectwise cofibrations and trivial cofibrations, so its ordinary right adjoint preserves trivial fibrations. The graph-factorisation proof of Ken Brown's lemma then shows that this right adjoint preserves weak equivalences between fibrant diagrams. Consequently every $u$ for which the ordinary $u_*$ exists is admissible under those uniform closed exact hypotheses. This answers the product, tower, and group-indexing questions in the scope where the required ordinary operations exist. The argument uses only finite products, the specified factorizations, lifting, and two-out-of-three; it does not infer the existence of missing infinite limits \citep[Lemmas~1.1.9 and~1.1.12]{Hovey1999}.

Grothendieck nevertheless gives a useful geometric proof for finite posets. A downward-closed cover reduces a limit to a pullback, and the fibrancy of the cospan makes that pullback homotopy invariant. The construction exposes the local geometry behind pointwise derived direct images. His subsequent noetherian induction is correct, but its proposed larger class has an unexpected limitation: a sober noetherian Alexandrov space associated to a poset is finite. The commentary proves the assertion directly. Pointwise computation for a general comma category also deserves its own hypothesis: restriction to that comma category must preserve the chosen fibrant diagrams. This holds for the finite-poset situation and in the usual cocomplete diagram-model setting, where the left adjoint to restriction along a discrete fibration is computed by fibrewise coproducts. Admissibility and the pointwise comparison are thus separately justified.

The derived-category examples illuminate a different boundary. For a finite shape, a diagram of bounded-below complexes has a common lower bound. For an infinite discrete shape, its components may be bounded below at degrees that tend to minus infinity. Accordingly, $\operatorname{Ch}^+(\mathcal A(I))$ and $\operatorname{Fun}(I^{\mathrm{op}},\operatorname{Ch}^+(\mathcal A))$ need not agree, and products need not stay in $D^+(\mathcal A)$. An explicit family of shifted copies of a field $k$ makes the obstruction visible. Unbounded resolutions answer the substantive existence question in the standard module and sheaf settings: $K$-injective resolutions allow right derived functors to be defined on the full unbounded derived category. Spaltenstein's work predates the present manuscript, and the general Grothendieck-category theorem is recorded precisely in the references. Beke's construction of the combinatorial injective model structure on unbounded complexes, combined with the diagram-model theorem, gives the all-small-category derivator for any Grothendieck abelian category \citep[Proposition~3.13]{Beke2000}. As a direct later algebraic treatment under the additional assumption of enough projectives, Alonso Tarr\'io, \'Alvarez D\'iaz, and Jerem\'ias L\'opez prove the all-small-category strong derivator theorem for a Grothendieck category with enough projectives \citep[Theorem~3.4, arXiv version~3]{AlonsoAlvarezJeremias2023}.

The modern comparison also separates pointedness from stability. Fibre and cofiber sequences are available in pointed homotopy theories without making their homotopy categories additive or their suspension functors invertible. The manuscript's deliberately broad ``triangul\'ee'' language must be read accordingly. Strong stable derivators give actual canonical triangulations, with exact diagram operations, as proved by \citet[Theorem~4.16 and Proposition~4.18]{Groth2013}. His Theorem~1.28 and Proposition~1.30 provide the combinatorial model-category construction of the full derivator. These results explain the successful modern form of the programme without replacing the original finite-limit framework by an unstated cocompleteness assumption.

Section~7 changes the technical language but keeps the construction problem in view. A cardinal filtration $(\Filt^\pi(E))$ consists of essentially small stages, bounded colimit-closure conditions, and presentations of every object by sufficiently directed diagrams in each stage. ``Grand devant $\pi$'' means that every subset of cardinality at most $\pi$ has an upper bound, hence means $\pi^+$-directed in the customary regular-cardinal convention. Membership in $\Filt^\pi(E)$ is a size condition; preservation of these colimits by $\operatorname{Hom}(A,-)$ is a presentability condition. The filtration axioms do not identify the two. SGA~4's comparison results explain why suitable cardinals permit one to use both, and why the extra hypothesis becomes visible in the discussion of canonical slice categories \citep[I, 9.12 and 9.16--9.18]{SGA4I1972}.

The adapted-presentation construction itself survives the audit well. Replace an indexing poset by its small sufficiently directed subposets and take the colimit over each piece. The new index admits the relevant small suprema, and the new diagram preserves them. The editors correctly observe that the first enlargement lemma is immediate under its stated hypotheses: one may simply adjoin an upper bound to the given small subset. The transfinite proof does something more useful when only the weaker directedness of the ambient poset is assumed. Its cardinal equation controls the number of newly imposed bounds. The commentary proves both versions and checks their different assumptions. The separate assertion that generalized continuum hypothesis always makes $(\pi')^\pi=\pi'$ is false at singular $\pi'$ with small cofinality; the displayed power-of-two case remains valid without that hypothesis.

The proposed equivalence of the three conditions on a part $M$ needs a more substantial correction. Condition $a_\pi$ describes $M$ as the objects obtainable by $\pi^+$-directed colimits from $M\cap\Filt^\pi(E)$. Condition $b_\pi$ asks, for every perfectly adapted presentation, that its colimit belong to $M$ if and only if its $M$-stages are cofinal. Condition $c_\pi$ combines strict generation, closure under sufficiently directed colimits, and ordinary filteredness of a canonical slice. These quantify over different data. The commentary gives two fully checked counterexamples in $E=\mathcal P(\omega_1)$ with $\pi=\aleph_0$. Taking $M$ to be the countable successor initial segments together with $\omega_1$ satisfies $a_\pi$, but a continuous diagram of countable limit initial segments violates $b_\pi$. Taking $M$ to be the finite subsets together with $\omega_1$ satisfies $c_\pi$ but violates $a_\pi$. Both examples use the ordinary cardinality filtration, whose stages have the required presentability. They therefore cannot be removed by adding only the compactness hypothesis mentioned in the manuscript.

What does remain is substantial. Grothendieck's proof of $b_\pi\Rightarrow a_\pi$ is completed without any presentability assumption, using a cofinal family of small subposets with greatest elements. With presentable small stages, the canonical slice construction gives the stronger directedness needed to prove colimit closure. Near page~94 there is a distinct size error: bounding the image of an indexing category does not bound that category itself. A filtered diagram of finite sets can have an uncountable colimit even though only countably many finite sets and maps occur as values. The repair bounds the indexing morphisms. Finally, the enlarged-cardinal idea is developed as a separate conditional result, with every object of $\Filt^\pi(E)$ assumed $\pi^+$-presentable and every object of $\Filt^{\pi^+}(E)$ assumed $\pi^{++}$-presentable, and with a complete alternating-chain proof. It is not presented as the missing end of Lemma~3. The powerset counterexamples and this conditional completion are marked as the commentators' work, with no claim of priority.

The accompanying registers preserve the distinction between a transcription mistake, a slip already present in the edition, and a mathematical claim requiring a proof or a correction. Every doubtful reading and every \texttt{[sic]} marker of the transcription is accounted for. The editorial indices, axiom list, and glossary have also been checked, including the missing cardinal bound in the glossary's definition of strict filtration. References report the portions actually verified. The purpose throughout is to let the reader follow the mathematical argument while retaining an exact account of what this fragment establishes, what later theory supplies, and what the present commentary proves separately.
\end{otherlanguage}
\dgrchunkend
